\newcommand{\DoNotLoadEpstopdf}{}
\documentclass[11pt]{article}
\usepackage[T1]{fontenc}
\usepackage{lmodern,amsmath,amssymb,amsthm,mathtools}
\usepackage[margin=1in]{geometry}
\usepackage{booktabs,array,longtable,enumitem,microtype}
\usepackage[hidelinks]{hyperref}
\usepackage{xurl}
\setlength{\emergencystretch}{2em}
\setlength{\LTpre}{.6\baselineskip}
\setlength{\LTpost}{.6\baselineskip}
\newtheorem{theorem}{Theorem}[section]
\newtheorem{proposition}[theorem]{Proposition}
\newtheorem{lemma}[theorem]{Lemma}
\newtheorem{corollary}[theorem]{Corollary}
\theoremstyle{definition}
\newtheorem{definition}[theorem]{Definition}
\newtheorem{assumption}[theorem]{Input}
\newtheorem*{example*}{Example}
\theoremstyle{remark}
\newtheorem{remark}[theorem]{Remark}
\numberwithin{equation}{section}
\newcommand{\HH}{\mathbb H}
\newcommand{\R}{\mathbb R}
\newcommand{\Q}{\mathbb Q}
\newcommand{\Z}{\mathbb Z}
\newcommand{\N}{\mathbb N}
\newcommand{\ip}[2]{\langle #1,#2\rangle}
\newcommand{\norm}[1]{\lVert #1\rVert}
\newcommand{\supp}{\operatorname{supp}}
\newcommand{\rank}{\operatorname{rank}}
\newcommand{\cone}{\operatorname{cone}}
\newcommand{\ph}{\varphi}
\newcommand{\Dfive}{\mathcal D_5}
\newcommand{\one}{\mathbf 1}
\title{Linear programming certificates for dimension bounds\\
of hyperbolic Coxeter polytopes}
\author{Alexander Kolpakov}
\date{Working manuscript, 18 November 2026}
\begin{document}
\maketitle

\begin{abstract}
We express three dimension arguments for hyperbolic Coxeter polytopes
as linear feasibility problems. Local geometric and arithmetic
theorems provide valid inequalities on complete configuration domains;
a common homogeneous formulation combines them with a strict global
budget. One certificate principle and one support-restriction argument
give the linear deduction and its dimension monotonicity. Finite
volume uses charges shared between face and vertex sums. Both compact
cases use supported energies of one original farthest vector and the
same capacity dual. Exact rational certificates recover the existing
bounds $57$, $17$ and $15$. We specify the inputs, explain their
geometric meaning and give the complete global deductions, without
repeating the local geometric proofs.
\end{abstract}

\tableofcontents

\section{Introduction}
\label{sec:introduction}

A dimension argument usually combines local restrictions on walls
and diagrams with a global count or energy budget. We put the second
step into a linear program. The variables retain the quantities
shared by different estimates; the dual determines how those
estimates may be combined. The local theorems remain explicit inputs.
This separates their geometric proofs from the global accounting and
allows the same assembly to be used with different input families.

Our applications recover the bounds in the companion geometric
treatment~\cite{BoundsArticle,BoundsSupplement}:
\begin{equation}\label{eq:three-bounds}
 \begin{aligned}
 n&\le57&&\text{in finite volume},\\
 n&\le17&&\text{in the compact case},\\
 n&\le15&&\text{in the compact arithmetic case}.
 \end{aligned}
\end{equation}
Guo and Zheng~\cite{GuoZheng2026} obtain the general finite-volume
bound $n\le81$. The present paper gives linear programming
certificates for the three bounds in~\eqref{eq:three-bounds}, using
the local inputs specified below.

Polytopes are finite-sided, convex and full-dimensional; hyperbolic
space has curvature $-1$. Arithmeticity refers to the reflection
group. Supplements FV, C and A mean, respectively, Parts I/IV, II/V
and III/VI of~\cite{BoundsSupplement}. We state their required local
results with their full domains rather than reproduce the geometric
proofs. Deriving those results from a smaller collection of Gram
conditions is a separate question.

\begin{theorem}[Linear programming recovery]
\label{thm:recovery}
The local estimates and complete configuration domains specified below
have the following exact rational certificates.
\begin{enumerate}[label=\textup{(\roman*)}]
\item The finite-volume incidence problem has a common charge
certificate for every integer dimension $n\ge58$. Its face payments
and vertex-budget inequalities contradict the strict incidence
excess in every such dimension.
\item The normalized compact energy problem on the complete
configuration family in dimensions $18,\ldots,29$ is infeasible.
A capacity certificate combines the local energy bounds into a
contradiction to the original vector's strict total budget.
\item The corresponding compact arithmetic problem in dimensions
$16,\ldots,29$ is infeasible. Its rank-$16$ and rank-$17$ arithmetic
domains and estimates are joined to the same compact
higher-dimensional domains and use the same capacity formulation.
\end{enumerate}
The two compact conclusions, together with the published compact
ceiling $29$, give the last two bounds in~\eqref{eq:three-bounds}.
The finite-volume certificate includes its own higher-dimensional
vertex and face estimates and gives the first bound.
\end{theorem}

The common mathematical structure is proved in
Theorem~\ref{thm:homogeneous-certificate}: an actual object supplies a
feasible point by normalization, and a linear functional rules it out.
Proposition~\ref{prop:threshold} proves monotonicity once, by restricting
one fixed dimension-labelled family. The applications supply different
valid constraints and different coverage proofs.

\begin{center}\small
\begin{tabular}{>{\raggedright\arraybackslash}p{.15\linewidth}
 >{\raggedright\arraybackslash}p{.41\linewidth}
 >{\raggedright\arraybackslash}p{.34\linewidth}}
\toprule
Case&Shared quantities and geometric budget&Linear certificate\\\midrule
Finite volume&Corner features counted in both face and vertex sums;
 the vertex allowances come from incidence&One charge pays the faces
 and remains within every vertex allowance\\
Compact&Orthogonal energies of one original vector, with strict
 total budget&Local multipliers respect each component's capacity
 and force the total up to the budget\\
Compact arithmetic&The same energies and budget, on arithmetically
 admitted domains&The same capacity dual uses the additional valid
 arithmetic rows and local exclusions\\\bottomrule
\end{tabular}
\end{center}

There are three logical roles to keep distinct. Geometric identities
and budgets hold on every actual object. Local theorems give further
inequalities or exclusions on stated domains. Certificate conditions
are restrictions on searched multipliers, such as face payments and
energy capacities; they are not additional hypotheses on a polytope.
Statements labelled \emph{Input} specify the first two roles. The LP
combination is proved here. In particular, the arithmetic pruning
used below is retained input, not a consequence newly derived by LP.

Uniformity is within each family: every component bound is used over
its full stated rank and spectator range, every ordinary dihedral
parameter keeps its complete range, and overlapping estimates share
the same original quantities. A formal domain may contain unrealizable
configurations. The necessary direction suffices for exclusion.
All terminating decimals denote exact rational numbers.

\section{Geometry, configurations and linear certificates}
\label{sec:framework}

\subsection{Wall conventions and spherical components}

We use a Lorentz space of signature $(-,+,\ldots,+)$, with hyperbolic
space the future unit hyperboloid. A wall has an outward spacelike
normal $u$ with $\ip{u}{u}=1$, and the chamber inequality is
$\ip{y}{u}\le0$. The Gram entry for incident Coxeter walls meeting at
angle $\pi/m$ is $-\cos(\pi/m)$. The Gram matrix at a finite vertex is
positive definite. Its connected components are mutually orthogonal
and span the Euclidean tangent space at that vertex.

We use the standard finite Coxeter classification. A canonical list
is
\[
 A_r\ (r\ge1),\quad B_r\ (r\ge3),\quad D_r\ (r\ge4),\quad
 E_6,E_7,E_8,F_4,H_3,H_4,\quad I_2(m)\ (m\ge4).
\]
Here $A_2=I_2(3)$, so the rank-two family has no missing label;
$B_2$ and $D_3$ are not counted as additional types. The \emph{order}
of a diagram is its number of nodes, and its \emph{rank} is the rank
of its Gram matrix. They agree for a spherical diagram. A connected
affine diagram with $r+1$ nodes has Gram rank $r$. The simplex factor
of an ideal vertex link that it determines has dimension $r$.
These conventions explain the ideal-link polynomials used later.

The geometric meaning of a local wall is part of its hypothesis.
When the source notation calls a diagram or label \emph{native}, it
means that of the specified actual walls, with their original node
numbering. This adjective never asserts that a projected or
replacement vector is another Coxeter wall.
In particular, an actual opposite wall is the facet reached on the
specified original edge. A replacement Gram matrix does not by itself
identify that wall. The local estimates cited below include the
appropriate actual-wall, finite-projection and farthest-vertex
arguments. Likewise, an arbitrary face is not assumed to inherit
the Coxeter or arithmetic hypotheses of the original polytope.

\subsection{The original vector in the compact case}

Let $P$ be compact, choose a strict interior point $z$, and choose a
vertex $p$ farthest from $z$. Put $H=-\ip{z}{p}=\cosh d(z,p)$.
The decomposition
\begin{equation}\label{eq:original-vector}
                   z=H(p-x),\qquad x\in p^\perp
\end{equation}
defines the original spatial vector $x$. Indeed,
$x=p-z/H$ is orthogonal to $p$, and $z^2=p^2=-1$ gives
\begin{equation}\label{eq:original-budget}
                         \norm{x}^2=1-H^{-2}<1.
\end{equation}
The space $p^\perp$ is Euclidean. If its orthogonal spherical
components are $W_1,\ldots,W_k$, define
\begin{equation}\label{eq:component-energies}
 e_j=\norm{\operatorname{proj}_{W_j}x}^2,\qquad
 e_j\ge0,\qquad \sum_{j=1}^ke_j=\norm{x}^2<1.
\end{equation}
An aggregate energy may group several of these mutually orthogonal
spaces. Every estimate added to this account must concern the same
vector $x$. An auxiliary point can prove an inequality for $x$;
its own energy is not an additional summand in~\eqref{eq:component-energies}.

\subsection{Complete domains and supported inequalities}

\begin{definition}\label{def:configuration-family}
A finite \emph{configuration family} at a dimension $n$ is a finite
collection of domains $\mathcal D_b$, each labelled by its actual
dimension $n_b=n$, such that every actual configuration under
consideration belongs to at least one $\mathcal D_b$. A domain records
all hypotheses needed by its local estimates. A complete symbolic
parameter range may be represented by one domain when the estimates
hold uniformly over that entire range.
\end{definition}

The domains need not be disjoint. Given an actual configuration one
may choose any matching domain, or fix a first-applicable rule. They
need not consist only of realizable configurations either. Extra
formal configurations make exclusion harder and do not compromise
the necessary direction.

A supported energy inequality consists of a set $S$ of original
orthogonal slots and a rational number $c\ge0$ such that
\begin{equation}\label{eq:supported-row}
                            \sum_{j\in S}e_j\ge c
\end{equation}
on its complete domain. If a theorem estimates a proper coordinate
subspace of one component, nonnegativity of the omitted energy gives
a bound for the whole component. Two bounds on overlapping subspaces
cannot then be added as though those subspaces were orthogonal.

A local theorem can instead prove that its entire domain is
impossible. We call this a \emph{local exclusion}. It is stated with
its full actual-wall or arithmetic hypotheses and will force the mass
of that domain to be zero. We keep it distinct from an energy bound.
A theorem excluding a complete component or a specified component
with a spectator is legitimate local input. The final exclusion of
all configurations at a dimension is the conclusion to be recovered.

\subsection{One homogeneous feasibility problem}

For reference we state the elementary certificate principle. Its use
in this paper always concerns a finite matrix, a justified finite
extended formulation, or a direct finite sum on an actual polytope.

\begin{lemma}\label{lem:cone-alternative}
Let $M$ and $d$ be finite rational data. Exactly one of the following
systems has a solution:
\begin{align}
 Mu=d,\quad u\ge0; &\label{eq:cone-primal}\\
 M^Ty\ge0,\quad d^Ty=-1. &\label{eq:cone-dual}
\end{align}
When either system is feasible, it has a rational solution.
\end{lemma}
\begin{proof}
Let $K$ be the cone generated by the columns of $M$. It is a closed
polyhedral cone. If $d\notin K$, separation gives a vector $y$ with
$y^Tk\ge0$ on $K$ and $y^Td<0$; rescaling gives
\eqref{eq:cone-dual}. If $d\in K$, its nonnegative column
coefficients give~\eqref{eq:cone-primal}. Both systems cannot be
feasible, since they would give
$-1=d^Ty=u^TM^Ty\ge0$.

For rationality, a nonempty polyhedron defined over $\Q$ has a
rational point: restrict to the affine hull of a minimal face,
solve its rational equality system, and approximate a relative
interior point by rational points in that affine space. All remaining
strict inequalities persist under a sufficiently small approximation.
Apply this to the two displayed rational systems.
\end{proof}

Free variables can be split into the difference of two nonnegative
variables. Linear inequalities can be converted to equalities by
slacks. Thus this lemma applies to the LPs below. Their more informative
duals, expressed as corner charges or energy capacities, are derived
explicitly rather than hidden in a large matrix transpose.

Here is a common formulation that also makes the dimension parameter
precise. Fix an index family $I$. Each index $i$ has an actual dimension
$n_i$, a polyhedral cone $C_i\subset\R^{k_i}$ of local variables,
a linear balance map $T_i:\R^{k_i}\longrightarrow W$, and a linear
functional $d_i:\R^{k_i}\longrightarrow\R$. The cones
contain homogeneous necessary inequalities. The map $T_i$ records
quantities that must balance between different local pieces, and
$d_i$ records their contribution to a strict global excess or deficit.
All these data are fixed before a dimension threshold is selected.

For an integer $N$ consider
\begin{equation}\label{eq:common-feasibility}
 \begin{gathered}
 u_i\in C_i\quad(n_i\ge N),\qquad
 \sum_{n_i\ge N}T_i u_i=0,\qquad
 \sum_{n_i\ge N}d_i(u_i)=1.
 \end{gathered}
\end{equation}
Only finitely many $u_i$ may be nonzero. If the family has infinitely
many dimensions, $W$ may have a separate block of balance coordinates
for each dimension. This prevents an incidence count at one dimension
from cancelling a count at another. An actual finite polytope always
uses finitely many variables.

\begin{theorem}[Homogeneous lift and certificate]
\label{thm:homogeneous-certificate}
Suppose every actual object of dimension $n\ge N$ in the covered
family supplies finitely supported local vectors $v_i\in C_i$,
supported on indices with $n_i=n$, with
\[
             \sum_iT_iv_i=0,\qquad D:=\sum_id_i(v_i)>0.
\]
Then it gives a feasible point of~\eqref{eq:common-feasibility}.
Conversely, a linear functional $\lambda$ on $W$ satisfying
\begin{equation}\label{eq:common-certificate}
       \lambda(T_iu)-d_i(u)\ge0
       \qquad(u\in C_i,\ n_i\ge N)
\end{equation}
certifies infeasibility and excludes those objects.
If $I$ is finite, $W$ is finite-dimensional and all the cones, maps
and functionals are rational, infeasibility is equivalent to the
existence of such a rational certificate.
\end{theorem}
\begin{proof}
Set $u_i=v_i/D$ for the lift. For a putative feasible point, summing
\eqref{eq:common-certificate} gives $0-1\ge0$, a contradiction.
For the finite converse, express each rational polyhedral cone by
finitely many rational generators and apply
Lemma~\ref{lem:cone-alternative} to their images
$(T_iu,d_i(u))$, with target $(0,1)$. Its separating functional has
the form $(\lambda,-1)$, exactly as in
\eqref{eq:common-certificate}.
\end{proof}

\begin{proposition}[Monotone thresholds]\label{prop:threshold}
In the fixed formulation~\eqref{eq:common-feasibility}, feasibility
at $N'\ge N$ implies feasibility at $N$. Thus infeasibility at $N$
implies infeasibility at every larger threshold. A certificate at
$N$ remains a certificate after restricting to $n_i\ge N'$.
\end{proposition}
\begin{proof}
Extend a feasible point by zeros on the removed indices. Every cone
contains zero, and both sums are unchanged. The certificate assertion
is restriction of the inequalities already proved.
\end{proof}

This argument uses neither dimension heredity nor a limiting duality
theorem. For an infinite family we use only the displayed certificate
on finite sums. A finite catalogue must be accompanied by a theorem
covering the dimensions outside it. The table in
Section~\ref{sec:thresholds} records that coverage for all three cases.
The following applications identify the cones, the balance and the
positive quantity being normalized.

\section{Finite volume: a charge on vertex--face incidences}
\label{fv:section}

Count each corner once at its face and once at its vertex. We seek a
single corner charge that pays the face demands while respecting the
vertex allowances. The resulting charge is a dual certificate for the
necessary incidence problem. Searching for it in a specified family is
a second LP; its optimum is a property of that family, rather than a
geometric assumption.

\subsection{Vertex figures and the incidence allowance}
\label{fv:incidence}

Fix an integer $n\ge58$. A \emph{corner} is an incidence
$\kappa=(v,F)$ of a vertex $v$ and a two-dimensional face $F$.
The number of sides of $F$ is denoted by $s(F)$, and
\begin{equation}\label{fv:defect}
 D_5(P)=\sum_{\dim F=2}(5-s(F)).
\end{equation}
A triangle contributes two, a quadrilateral contributes one, and a
face with at least five sides contributes a nonpositive number.

We use the standard vertex-figure description of a finite-volume
Coxeter polytope. At a finite vertex the figure is a simplex. At an
ideal vertex a horospherical figure is a product
\begin{equation}\label{fv:link}
 L_v=\prod_{i=1}^k\Delta^{r_i},\qquad
 r_i\ge1,\qquad \sum_i r_i=n-1.
\end{equation}
Here $r_i$ is the dimension of a simplex factor, not the number of
nodes of its affine diagram. That diagram has $r_i+1$ nodes and Gram
rank $r_i$. These vertex-figure statements, including their applicability
to nonsimple ideal vertices, are retained geometric inputs from
\cite{BoundsArticle,BoundsSupplement}.

Put $r=n-1$. Write the $h$-polynomial of the figure as
\begin{equation}\label{fv:hpolynomial}
 h_v(t)=\sum_{j=0}^r h_{v,j}t^j,
 \qquad
 h_v(t)=\prod_i(1+t+\cdots+t^{r_i})
 \quad\hbox{at an ideal vertex}.
\end{equation}
At a finite vertex $h_{v,j}=1$ for every $0\le j\le r$.
Define
\begin{equation}\label{fv:allowance-general}
 V_v=\sum_{j=0}^r h_{v,j},\qquad
 G_v=\sum_{j=0}^r h_{v,j}\min(j,r-j),\qquad
 b_n(v)=\frac52G_v-\frac r2V_v.
\end{equation}
The quantity $V_v$ is the number of vertices of the vertex figure.
It equals $n$ at a finite vertex and $\prod_i(r_i+1)$ at an ideal
vertex. Every vertex of the figure has degree $r$, so its number of
edges is $E_v=rV_v/2$. These edges correspond to corners at $v$.

For a finite vertex, the sum in \eqref{fv:allowance-general} gives
\begin{equation}\label{fv:finite-allowance}
 \beta_n=b_n(v)=
 \begin{cases}
 n(n-6)/8,&n\text{ even},\\
 (n-1)(n-5)/8,&n\text{ odd}.
 \end{cases}
\end{equation}
At an ideal vertex we write
\begin{equation}\label{fv:ideal-allowance}
 K_n(v)=\frac{5\sum_jh_{v,j}\min(j,n-1-j)-(n-1)V_v}{2}.
\end{equation}
Thus $K_n(v)=b_n(v)$; the separate letter only distinguishes the
ideal allowance from the constant finite allowance $\beta_n$.

The incidence input needed below has the following precise form.

\begin{assumption}[Incidence excess]\label{fv:incidence-excess}
Let $P\subset\HH^n$ be a finite-volume Coxeter polytope, with the
vertex figures described above. The incidence theorem of
\cite{BoundsSupplement} gives
\begin{equation}\label{fv:incidence-excess-eq}
 D_5(P)\ge\sum_v b_n(v)+\Delta_n,
 \qquad \Delta_n\ge5(n-1)>0.
\end{equation}
\end{assumption}

This imports the height-index argument, including its strict excess,
as one geometric input. The LP need only prove
$D_5(P)\le\sum_vb_n(v)$; it need not produce a positive charge margin.

\subsection{The incidence feasibility problem and its charge certificate}
\label{fv:incidence-primal}

We first give the formulation without choosing particular corner
features. Fix a vector of rational features $\Psi(\kappa)\in\Q^d$.
The same feature vector must be used when a corner is viewed from its
face and when it is viewed from its vertex. A vertex type $V$ includes
its complete decorated vertex diagram or product-link data. A small-face
type $Q$ includes a triangle or quadrilateral together with its coherent
corner descriptions. Set
\begin{equation}\label{fv:feature-sums}
 A(V)=\sum_{\kappa\text{ at }V}\Psi(\kappa),\qquad
 A(Q)=\sum_{\kappa\text{ on }Q}\Psi(\kappa),\qquad
 d(Q)=5-s(Q).
\end{equation}
All domains used here are necessary domains: every actual vertex and
small face must have an allowed description. An allowed description
need not extend to an actual polytope.

Introduce nonnegative variables $x_V$ for vertex types, $y_Q$ for
small-face types, and $z_\kappa$ for allowed individual corners. The
last variables account for corners belonging to faces with at least
five sides. The incidence feasibility problem is
\begin{align}
 \sum_VA(V)x_V-\sum_QA(Q)y_Q-
       \sum_\kappa\Psi(\kappa)z_\kappa&=0,
       \label{fv:primal-balance}\\
 \sum_Qd(Q)y_Q-\sum_Vb_n(V)x_V&=1.
       \label{fv:primal-normalization}
\end{align}
If a domain is infinite, the variables may be read as finitely
supported nonnegative measures. Actual polytopes use only finitely
many incidences, which is all that the argument needs.

\begin{lemma}\label{fv:actual-lift}
An actual finite-volume Coxeter polytope gives a feasible point of
\eqref{fv:primal-balance}--\eqref{fv:primal-normalization} whenever
the allowed domains contain all its local descriptions.
\end{lemma}

\begin{proof}
Let $f_3$ and $f_4$ be the numbers of triangular and quadrilateral
two-faces, and put
\[
 L=2f_3+f_4-\sum_vb_n(v).
\]
Since the other two-faces contribute nonpositively to $D_5(P)$,
Input~\ref{fv:incidence-excess} gives
$L\ge\Delta_n>0$. Divide each actual vertex count, small-face count
and residual-corner count by $L$. Counting the same corners by vertices
and by faces gives \eqref{fv:primal-balance}. The definition of $L$
gives \eqref{fv:primal-normalization}.
\end{proof}

The residual variables in \eqref{fv:primal-balance} are important.
Their contribution has not disappeared merely because their faces have
nonpositive demand. Feature conservation still counts those corners.

\begin{proposition}[Charge certificate]\label{fv:charge-certificate}
Suppose a coefficient vector $\theta\in\R^d$ satisfies
\begin{align}
 b_n(V)-\theta\cdot A(V)&\ge0 &&\text{for every allowed vertex type},
    \label{fv:vertex-certificate}\\
 \theta\cdot A(Q)-d(Q)&\ge0 &&\text{for every allowed small face},
    \label{fv:face-certificate}\\
 \theta\cdot\Psi(\kappa)&\ge0 &&\text{for every allowed corner}.
    \label{fv:corner-certificate}
\end{align}
Then the incidence feasibility problem is infeasible.
\end{proposition}

\begin{proof}
Use Theorem~\ref{thm:homogeneous-certificate}, reversing the sign of
\eqref{fv:primal-balance}. The vertex, small-face and residual-corner
columns $(T,d)$ are respectively $(-A(V),-b_n(V))$,
$(A(Q),d(Q))$ and $(\Psi(\kappa),0)$; take $\lambda=\theta$.
The three displayed conditions are exactly its certificate inequalities.
\end{proof}

Equivalently, $w(\kappa)=\theta\cdot\Psi(\kappa)$ pays every triangle
and quadrilateral and stays within every vertex allowance. Faces with
at least five sides are already paid by nonnegative corner charges.
Summing gives $D_5(P)\le\sum_vb_n(v)$, contradicting incidence.
This is a direct weak-duality argument. It makes no use of
infinite-dimensional strong duality or a realization theorem for the
allowed local descriptions.

\subsection{The corner features used in the recovered certificate}
\label{fv:features}

We now specify the finite-dimensional family used in the actual
certificate. Its basic feature vector is
$\phi(\kappa)\in\Q^{26}$. A feature is a function of the decorated
corner, whereas its coefficient is an optimization variable.

At a finite vertex, the spherical diagram is a disjoint union of
connected finite diagrams. Its two marks are the two nodes determining
the corner. At an ideal vertex a corner corresponds to an edge of
the product \eqref{fv:link}: one simplex factor varies along an edge,
and a vertex is selected in every other factor. The varying affine
diagram has two marked nodes. Each other factor contributes the finite
diagram obtained by deleting the node selected by its vertex. These
finite residual diagrams are retained in the corner description.

We number $A_r$ along its path by $0,\ldots,r-1$. We number $D_r$
along the spine $0,\ldots,r-3$, with leaves $r-2,r-1$ attached to
$r-3$. The remaining finite and affine diagrams use the fixed
numberings of the rational tables in \cite{BoundsSupplement}.
One numbering is selected for a complete vertex diagram and used at
all its corners. A relabelling must carry the numbering with it; one
cannot choose a new favourable numbering separately at every corner.

Set $\phi_0=1$. Features $\phi_1,\phi_2,\phi_3$ are respectively
the prescribed ordinary pair weight for two marks in the same finite
component, the affine pair weight in the varying ideal factor, and
the ordinary weight for marks in different finite component copies;
each is zero outside its stated case. Their fixed rational tables,
including exceptional small ranks, are the intrinsic corner tables of
\cite{BoundsSupplement}. We use their final edition: in particular,
the ordinary $B_{13}$ endpoint entry is zero.

The next fifteen features are indicators. Each is one when the two
marks lie in a copy of the displayed component at the displayed
unordered pair, and when the final column holds. Otherwise it is zero.
A component mentioned as a witness is a connected component of the
complete corner diagram disjoint from the marked component. Put
$M_{15}=m_{A_{15}}+m_{D_{15}}$, the number of such finite components.
\begin{center}\small
\begin{tabular}{cllp{.30\textwidth}}
\toprule
$j$&Marked type&Marked pair&Witness condition\\\midrule
4&$A_{13}$&$\{0,12\}$&An $A_1$ occurs.\\
5&$A_{14}$&$\{0,13\}$&An $A_1$ occurs.\\
6,7&$A_{15}$&$\{0,13\},\{1,14\}$&An $A_1$ occurs.\\
8,9&$D_{15}$&$\{1,13\},\{1,14\}$&An $A_1$ occurs.\\
10,11&$A_{16}$&$\{1,12\},\{3,14\}$&An $A_1$ occurs.\\
12,13&$D_{16}$&$\{0,14\},\{0,15\}$&An $A_1$ occurs.\\
14,15&$D_{14}$&$\{0,12\},\{0,13\}$&An $A_1,A_2$ or $A_3$ occurs.\\
16,17&$D_{13}$&$\{0,11\},\{0,12\}$&$M_{15}\ge3$.\\
18&$D_{28}$&$\{0,12\}$&$M_{15}\ge2$.\\
\bottomrule
\end{tabular}
\end{center}
In a row with two feature indices, the two marked pairs correspond
in the displayed order.
The final seven indicators concern different finite component copies:
\begin{center}\small
\begin{tabular}{cp{.48\textwidth}p{.31\textwidth}}
\toprule
$j$&Marked components and allowed nodes&Witness condition\\\midrule
19&$A_4:\{0,3\}$ and $A_{12}:\{1,10\}$&An $A_1$ occurs.\\
20&$A_2:\{0,1\}$ and $A_4:\{0,3\}$&An $A_1$ and an $A_{10}$ occur.\\
21&$A_2:\{0,1\}$ and $D_5:\{0\}$&An $A_1$ and an $A_9$ occur.\\
22&$A_1:\{0\}$ and $D_5:\{0\}$&An $A_{11}$ occurs, or both an $A_9$ and an $A_2$ occur.\\
23&$A_1:\{0\}$ and $D_5:\{0\}$&An $A_9$ and an $A_2$ occur.\\
24&$A_{12}:\{0,11\}$ and $D_{16}:\{0\}$&$M_{15}\ge2$.\\
25&Two distinct $A_2$ copies, with either node in each&An $A_{10}$ and at least two $A_1$ copies occur.\\
\bottomrule
\end{tabular}
\end{center}
These are indicators of unordered marked pairs, not oriented-pair
counts. Features 22 and 23 deliberately overlap. Feature 25 vanishes
when the two marks lie in one $A_2$ copy. The raw features all lie in
$[0,1]$; this does not imply nonnegativity of an arbitrary linear
combination of them.

There is one additional ideal-corner indicator, denoted $\chi_B$.
Write $\Gamma[m]$ for an affine diagram with $m$ nodes and rank $m-1$.
The indicator is one precisely when the varying factor is $B[13]$,
the marked pair is $\{10,11\}$ or $\{10,12\}$, and the complete
residual finite diagram contains at least three components of types
$A_{15}$ or $D_{15}$. It is zero otherwise. Thus $B[13]$ denotes
standard type $\widetilde B_{12}$, not a finite rank-13 component.
The residual diagrams are formed by the other factors' one-node
deletions. One does not first delete the two marked nodes and count
newly separated pieces in the varying factor.

Let $\epsilon_n=[n=58]$. Use free vectors
$a^{58},a^{\rm tail}\in\R^{26}$ and parameters $\lambda,\rho\ge0$.
The searched charge is
\begin{equation}\label{fv:piecewise-charge}
 w(\kappa)=
 \bigl(\epsilon_na^{58}+(1-\epsilon_n)a^{\rm tail}\bigr)
       \cdot\phi(\kappa)+\epsilon_n\lambda\chi_B(\kappa).
\end{equation}
Its 53 conserved features are
\begin{equation}\label{fv:conserved-features}
 \Psi(\kappa)=
 (\epsilon_n\phi(\kappa),(1-\epsilon_n)\phi(\kappa),
                     \epsilon_n\chi_B(\kappa)).
\end{equation}
The margin $\rho$ is a separate certificate variable, giving 54
shared variables in total. It is not another corner feature.
The coefficient vectors are free: some entries in the successful
certificate are negative. Nonnegativity is imposed on the resulting
charge on admitted corners, rather than on every coefficient.

The finite table construction used in the checked endpoint network
has a declared ambient range through rank 58. No assertion about
larger dimensions is obtained by running that constructor outside
its range. The meaning and validity of \eqref{fv:piecewise-charge}
above the endpoint come from the separately retained all-rank face
theorems and analytic vertex estimates described below.

\subsection{The geometric face inputs and their linear conclusions}
\label{fv:face-inputs}

The small-face theorems determine which marked components and original
supporting walls can occur together, and which estimates can pay one
face. We retain their complete domains as the following LP interface.

\begin{assumption}[Retained local face theorems]\label{fv:face-assumption}
Use the complete corner, triangle and quadrilateral coverage of
\cite{BoundsSupplement}, together with the common-coefficient
transfers recorded in \cite{LPRepository}. Their coefficient conditions
form the declared local linear system $\mathcal L$. It contains the
corner and triangle rows on both dimension branches; the guarded
ideal and finite $D_{13}$ conditions; the fork and additional
rank-16 conditions; and the complete retained quadrilateral
coefficient conditions, including both branch-specific $A_{13}$ rows.
For any feasible shared coefficients, with the stated auxiliary
certificate variables, these theorems imply
\eqref{fv:face-certificate}--\eqref{fv:corner-certificate} for every
actual corner, triangle and quadrilateral in every $n\ge58$.
When a local transfer uses equalities forced by the endpoint vertex
budgets, its hypotheses include the complete network conditions of
Section~\ref{fv:finite-network}; those equalities are not presumed as
independent geometric axioms.
\end{assumption}

The word ``conditions'' in this input refers to restrictions
on searched charge coefficients. It does not assert that a polytope
has to satisfy numerical inequalities in variables $a^{58}$ and
$a^{\rm tail}$. Those variables describe a proof we are constructing.
Nor are the conditions claimed necessary for every possible valid
charge. For example, the retained $A_{13}$ argument uses
\begin{equation}\label{fv:A13-rows}
 a^{58}_0+a^{58}_4\ge1,
 \qquad a^{\rm tail}_0+a^{\rm tail}_4\ge1.
\end{equation}
These rows are part of this specified sufficient certificate family.

The coverage retains actual wall identities, ordered marks, complete
maps into the fixed numbered diagrams, residual choices at ideal links,
and every original wall used by an auxiliary construction. These are
the applicability conditions of the imported theorem; an inverse-Gram
calculation alone does not supply them. Face payments from different
alternatives are never added as though they belonged to disjoint corners.

Two examples show the actual linear form. For either branch vector
$z=a^{58}$ or $a^{\rm tail}$, the retained conditions include
\begin{align}
 4z_0+z_3/16+z_{22}+z_{23}&\ge1,
       \label{fv:example-face-P}\\
 4z_0+z_3/2+z_{21}&\ge1.
       \label{fv:example-face-M}
\end{align}
The geometric domains on which these rows pay the face are part
of their respective local theorems. The coefficients are not inferred
from the appearance of a matching type name elsewhere in a face.
The electronic local matrix gives every retained row and its source
association. Its geometric validity is the content of the cited local
results, not a consequence of the matrix having rational entries.

\subsection{Finite vertices: an exact vector-cost network}
\label{fv:finite-network}

At a finite vertex the marks range over all unordered pairs of
nodes in a spherical forest. A formal forest has nonnegative integer
component multiplicities $m_T$ satisfying
\begin{equation}\label{fv:forest-rank}
                  \sum_T r_Tm_T=58,
\end{equation}
where $r_T$ is the rank of type $T$. We use the finite collection of
weight types of the rational catalogue. A symbolic dihedral family
can be represented by one weight type only when its complete feature
table is the same throughout the stated parameter range and that
range is covered by the local theorem. A finite sample of labels is
not used in place of an unbounded family.

Let
\[
 \Phi(m)=\sum_{\kappa\text{ at the formal vertex}}\phi(\kappa).
\]
Every actual finite vertex supplies an admitted formal forest. The
converse is not required. We seek the family of inequalities
\begin{equation}\label{fv:finite-budget-endpoint}
                   a^{58}\cdot\Phi(m)+\rho\le377
                  \quad\hbox{for every }m\text{ satisfying
                  \eqref{fv:forest-rank}}.
\end{equation}
The right side is the incidence allowance
$\beta_{58}=58(58-6)/8$, rather than an assumed optimized vertex cost.

Here is the structure that makes this large family a network LP.
For one component $T$, let $S_T$ be the sum of its ordinary
same-component entries; for different component types $T,U$, let
$B_{TU}$ be the sum of their ordinary cross entries. Using $u_j$
for the $j$th coordinate vector of $\Q^{26}$ gives
\begin{equation}\label{fv:first-aggregates}
 \Phi_0=\binom{58}{2},\qquad
 \Phi_1=\sum_Tm_TS_T,\qquad
 \Phi_2=0,\qquad
 \Phi_3=\sum_{T<U}m_Tm_UB_{TU}.
\end{equation}
Features 4--18 sum to the marked-type multiplicity times their witness
indicator. The remaining coordinates follow by counting the unordered
mark pairs in the table: for example,
$\Phi_{19}=4m_{A_4}m_{A_{12}}[m_{A_1}>0]$ and
$\Phi_{25}=4\binom{m_{A_2}}2[m_{A_{10}}>0,m_{A_1}\ge2]$.
Brackets denote Boolean indicators; the factors four count node choices,
not orientations.

Fix the multiplicities $h_H=m_H$ of the types in
\begin{equation}\label{fv:high-set}
 \mathcal H=\{A_9,\ldots,A_{17},B_9,\ldots,B_{12},
                    D_9,\ldots,D_{16},D_{28}\}.
\end{equation}
Also fix $a=m_{A_1}$ and $b=m_{A_2}$; call the resulting tuple an
interface $I$. Let
\[
 M_h=58-\sum_Hr_Hh_H,\qquad R_I=M_h-a-2b.
\]
All nonzero ordinary cross terms have exactly one endpoint in
$\mathcal H$. After an interface is fixed, the remaining vector cost
is additive in the residual component counts, except for one
$A_3$-presence condition. More precisely, let $\mathcal R$ be the
remaining type alphabet after removing $\mathcal H,A_1,A_2$, and put
$\mathcal S=\mathcal R\setminus\{A_3\}$. Define
\[
 p_T^h=S_Tu_1+\Bigl(\sum_Hh_HB_{TH}\Bigr)u_3.
\]
For residual types other than $A_4,D_5$ use $p_{I,T}=p_T^h$.
The two exceptions have the complete vector costs
\begin{align}
 p_{I,A_4}={}&p_{A_4}^h+4h_{A_{12}}[a>0]u_{19}
                    +4b[a>0,h_{A_{10}}>0]u_{20},\notag\\
 p_{I,D_5}={}&p_{D_5}^h+2b[a>0,h_{A_9}>0]u_{21}\notag\\
 &+a[h_{A_{11}}>0\text{ or }(h_{A_9}>0,b>0)]u_{22}
                    +a[h_{A_9}>0,b>0]u_{23}.
                    \label{fv:residual-costs}
\end{align}
Let $K_I$ contain $\binom{58}{2}u_0$ and all terms of the feature
sum involving only the fixed counts $h,a,b$, including the fixed
same-component indicators. Put
\[
              \gamma_I=h_{D_{14}}[a=b=0](u_{14}+u_{15}).
\]
Then, with $k=m_{A_3}$,
\begin{equation}\label{fv:vector-cost-identity}
 \Phi(m)=K_I+kp_{I,A_3}+[k>0]\gamma_I
                         +\sum_{T\in\mathcal S}m_Tp_{I,T}.
\end{equation}
This identity follows directly from the feature definitions: the only
unfixed witness is $A_3$ in the $D_{14}$ indicators, the unfixed cross
terms are linear in $A_4,D_5$, and the others are fixed by $I$.

An entry for $(I,k)$ has vector cost
$K_I+kp_{I,A_3}+[k>0]\gamma_I$ and residual rank $R_I-3k$.
A transition removes a type $T\in\mathcal S$, reduces rank by
$r_T$, and contributes vector cost $p_{I,T}$. Only rank-feasible
interfaces, entries and reachable residual states occur. Rank strictly
decreases along every transition, so the graph is acyclic.

\begin{lemma}[Exact network formulation]\label{fv:network-exact}
Every admitted formal forest has a path whose vector cost is
$\Phi(m)$, and every path gives admitted counts of total rank 58.
Consequently the maximum path charge equals
$\max_m a^{58}\cdot\Phi(m)$ for every signed vector $a^{58}$.
\end{lemma}

\begin{proof}
The interface, entry and residual transitions select exactly the
nonnegative counts in \eqref{fv:forest-rank};
\eqref{fv:vector-cost-identity} gives the path cost. Different removal
orders can repeat a forest but cannot change its cost. A unit flow on
this acyclic graph is a convex combination of complete paths, so its
maximum linear cost is the maximum path cost.
\end{proof}

For a fixed high assignment, interfaces may share a residual network
only if their complete vectors $p_{I,A_4},p_{I,D_5}$ agree. The other
residual costs are already fixed by $h$. Such sharing retains every
entry separately. Equality of scalar costs at one previously selected
charge would not justify sharing while the charge coefficients vary.

For a network fragment $S$, write $v_{S,g,\ell}$ for free potentials
at reachable residual rank $\ell$ in sharing group $g$. Let $U_S$ be a
free fragment upper bound. The rows are
\begin{align}
 v_{S,g,0}&=0,\notag\\
 v_{S,g,\ell}-v_{S,g,\ell-r_T}
                  &\ge a^{58}\cdot p_{g,T},\notag\\
 U_S-v_{S,g(I),R_I-3k}
                  &\ge a^{58}\cdot(K_I+kp_{I,A_3}+[k>0]\gamma_I),
                  \label{fv:network-lp}\\
 U_S+\rho&\le377.\notag
\end{align}
All fragments use the same $a^{58},\rho$ and retain every residual
type; only potentials are private. Telescoping gives every path bound.
Conversely, maximum residual-completion costs give feasible potentials.
Eliminating them therefore gives exactly
\eqref{fv:finite-budget-endpoint}, even for signed coefficients.
Neither integral optimization variables nor nonnegative potentials
are needed. Completeness of the checked graph is a finite-data input;
the lemma explains why its encoding is an LP.

\subsection{Ideal vertices and the analytic tail}
\label{fv:ideal-block}

Let an ideal link have factor ranks
$r_1,\ldots,r_k$, total rank $R=n-1$ and $V=\prod_i(r_i+1)$ vertices.
For a connected affine factor $\Gamma$ of rank $r$, let $Q_\Gamma$
be the sum of its ordinary marked-pair weights and put
$q_\Gamma=Q_\Gamma/(r+1)$. Define $C=\sum_iq_{\Gamma_i}$.
The normalized ordinary charge for coefficient values $(a_0,a_2)$ is
\begin{equation}\label{fv:ideal-charge}
              \frac{C(v)}{V}=\frac R2a_0+Ca_2
\end{equation}
before adding the guarded indicator. This follows because a marked
pair in one factor occurs once for each choice of a vertex in every
other factor. Only raw features 0 and 2 occur at ideal corners.

Let $J$ be the sum of independent uniforms on
$\{0,\ldots,r_i\}$ and put $Y=2J-R$. The coefficient distribution
of $h_v$ is the distribution of $J$ multiplied by $V$. Symmetry gives
\begin{equation}\label{fv:normalized-ideal}
 \frac{K_n(v)}V=\frac{3R}{4}-\frac54\mathbb E|Y|,
 \qquad
 S_a=\frac{3R}{4}-\frac54\mathbb E|Y|-\frac R2a_0-Ca_2.
\end{equation}
The desired vertex inequality is $S_a\ge\rho$, so the raw ideal
margin is $V\rho$, not an unscaled copy of $\rho$.

\paragraph{The rank-57 block.}
For $n=58$, introduce free factor-cost variables $k_r$ for
$1\le r\le57$. For every connected affine family impose
\begin{equation}\label{fv:ideal-factor-links}
                 k_r\ge\frac r2a^{58}_0+q_\Gamma a^{58}_2.
\end{equation}
The complete family ranges are affine $A$ at ranks 1--57,
$C$ at 2--57, $B$ at 3--57, $D$ at 4--57, and the exceptional
$E_6,E_7,E_8,F_4,G_2$. There are 227 individual-family rows.
They all remain, even when $a^{58}_2$ is negative.

The ideal input supplies 354 retained symmetric payoff vectors, with
fixed ancestry and domination indices covering every factor addition.
Domination compares all central-block averages with those of one
retained vector; its index cannot vary with the coordinate. The indices
are data, while the payoff coordinates below are LP variables. The
underlying operation on an array is uniform convolution minus its
factor cost:
\[
 f'_j=\frac1{r+1}\sum_{\ell=0}^r f_{j+\ell}-k_r.
\]
Its initial array is $g_j=(5\min(j,57-j)-57)/2$; its final scalar
is the normalized allowance minus the factor costs. At accumulated
rank $R$, put $L=57-R$ and let
$B_{R,i,a}$ be the payoff average on the
central interval $a,\ldots,L-a$, for
$0\le a\le\lfloor L/2\rfloor$. Let $I_R$ be the retained index
set at level $R$. If a rank-$r$ factor is added to parent
$j\in I_{R-r}$, put $A=L-2a+1$ and
\begin{align}
 \omega_{L,a,r,\ell}&=\frac{A+r-2\ell}{A(r+1)},
       &&0\le\ell<\min(A,r+1),\notag\\
 T_{R,i,a}^{r,j}&=B_{R,i,a}+k_r-
       \sum_{\ell=0}^{\min(A,r+1)-1}
           \omega_{L,a,r,\ell}B_{R-r,j,a+\ell}.
       \label{fv:ideal-average-transform}
\end{align}
These positive coefficients sum to one and express the central average
of the convolution as an average of nested central intervals. The rows
are:
\begin{enumerate}[label=\textup{(\roman*)},nosep]
\item the 29 initial equalities
\[
 B_{0,0,a}=\frac1{58-2a}\sum_{j=a}^{57-a}
                  \frac{5\min(j,57-j)-57}{2},
 \qquad0\le a\le28;
\]
\item $T_{R,i,a}^{r,j}=0$ for the fixed ancestry $(r,j)$ of
 each retained $i\in I_R$;
\item $T_{R,D(R,r,j),a}^{r,j}\le0$ for every candidate, using
 its fixed retained dominator $D(R,r,j)$;
\item $\tau\le B_{57,i,0}$ for every final retained index, and
 $\rho\le\tau$, with $\tau$ free.
\end{enumerate}
The fixed index sets and maps are specified in the retained electronic
ideal data of \cite{LPRepository}. With the family links this block has
103,265 rows and 5,754 private variables. Positive convolution propagates
the retained lower bounds through every factor addition, so every
rank-57 profile has surplus at least $\tau\ge\rho$. These fixed
domination choices define a sufficient certificate family, rather than
all possible ideal-vertex laws.

\paragraph{The guarded ideal feature.}
The additive block just described remains in the LP. The additional
feature $\chi_B$ is handled separately; replacing its effect by an
unconditional change of the rank-12 factor cost would be incorrect.
If $\chi_B$ occurs, the varying $B[13]$ factor has rank 12 and
the residual finite rank is 45. Three rank-15 components exhaust that
rank. The complete one-node-deletion classification leaves exactly
\begin{align}
 \mathrm{I}:&\quad B[13]\times\Gamma_1[16]
                   \times\Gamma_2[16]\times\Gamma_3[16],\notag\\
 \mathrm{II}:&\quad B[13]\times D[31]\times\Gamma[16],
 \qquad\Gamma_i,\Gamma\in\{A,B,D\}.
 \label{fv:guarded-links}
\end{align}
Let $a_\chi=V^{-1}\sum_\kappa\chi_B(\kappa)$. The additional rows are
\begin{equation}\label{fv:guarded-rows}
  \frac{57}{2}a^{58}_0+C a^{58}_2+a_\chi\lambda+\rho\le K.
\end{equation}
All thirteen rows have the following compact specification. Put
\[
 (c_A,c_B,c_D)=(5,25/4,411/64),\qquad
 (d_A,d_B,d_D)=(1,1/16,1/4).
\]
For every unordered triple $\Gamma_1,\Gamma_2,\Gamma_3\in\{A,B,D\}$
in type I, and every $\Gamma\in\{A,B,D\}$ in type II, use
\begin{equation}\label{fv:guarded-data}
\begin{array}{c|ccc}
 &C&a_\chi&K\\\hline
 \mathrm{I}&11/2+\sum_{i=1}^3c_{\Gamma_i}
       &(2/13)\prod_{i=1}^3d_{\Gamma_i}&331643/13312\\
 \mathrm{II}&753/62+c_\Gamma&(2/403)d_\Gamma&16549/806.
\end{array}
\end{equation}
Here $d_\Gamma$ is the proportion of one-node deletions giving the
required residual component; the deletion counts are $16,1,4$.
The values of $K$ are \eqref{fv:normalized-ideal} for ranks
$(12,15,15,15)$ and $(12,30,15)$. Thus the specification retains
deletion multiplicities, not just the affine factor types.

\paragraph{The residual dimension-59 sector.}
Append all 63 connected-family links of ranks at most 16 for the
tail coefficients,
\begin{equation}\label{fv:tail-factor-links}
             k_r\ge\frac r2a^{\rm tail}_0
                                      +q_\Gamma a^{\rm tail}_2,
\end{equation}
using the \emph{same} costs $k_r$ as above. Put $k_0=0$ and introduce
$\mu\ge0$. The fifteen descent rows are
\begin{equation}\label{fv:descent}
 k_r-k_{r-1}+\mu\le h_r,\qquad 1\le r\le15,
\end{equation}
where
\[
 (h_1,\ldots,h_{15})=
 \left(\frac34,\frac{47}{66},\frac{89}{132},\frac7{11},
 \frac{79}{132},\frac{37}{66},\frac{23}{44},
 \frac{2213}{4374},\frac{59}{132},\frac9{22},
 \frac{67}{156},\frac{71}{168},\frac5{12},
 \frac{79}{192},\frac{83}{204}\right).
\]
The retained descent theorem states that a rank-58 product with
largest factor rank between 10 and 16 has a factor whose rank can be
reduced by one so that its envelope surplus increases on restoring
that factor by at least $\mu$. The reduced product has rank 57 and
is covered by the full preceding block. This proves the required
dimension-59 ideal budget in that sector. Independent copies of the
costs in the two blocks would not justify this argument.

\paragraph{The unbounded ideal tail.}
For every remaining dimension-59 ideal link, and every ideal link
in dimension at least 60, the retained analytic theorem gives surplus
at least $\varepsilon=1/5307$ for baseline coefficients $(0,1)$.
For an arbitrary tail law a sufficient interval test reduces exactly
to four affine rows:
\begin{align}
 a^{\rm tail}_0&\le0,&
 a^{\rm tail}_0+a^{\rm tail}_2&\le1,\notag\\
 \rho+29a^{\rm tail}_0&\le\varepsilon,&
 \rho+29a^{\rm tail}_0+29a^{\rm tail}_2&\le29+\varepsilon.
 \label{fv:analytic-ideal-rows}
\end{align}
To see the reduction, write $t=1/n$,
$u_0=a^{\rm tail}_0$, $u_1=a^{\rm tail}_0+a^{\rm tail}_2-1$ and
\[
 P_i(t)=(\varepsilon-\rho)t-\frac{1-t}{2}u_i.
\]
Each $P_i$ is affine, and \eqref{fv:analytic-ideal-rows} is exactly
nonnegativity at the two endpoints of $[0,1/59]$. The analytic input
also gives $0\le C\le R/2$. If $w=2C/R$, then
\[
 t(S_a-\rho)=t(S_{(0,1)}-\varepsilon)
                         +(1-w)P_0(t)+wP_1(t)\ge0.
\]
Thus these four rows prove the desired budget on the entire stated
tail. Their validity uses the analytic surplus theorem; it is not
inferred from checking finitely many dimensions.

\paragraph{Finite vertices above the endpoint.}
Let $q\in\Q^{26}$ be the vector displayed in
\eqref{fv:successful-law} below. Introduce $0\le\xi\le1$ and impose
\begin{equation}\label{fv:finite-tail-rows}
 a^{\rm tail}_j\le q_j+2\xi[\,j\in\{16,17\}\,]
       \quad(j\ne2),\qquad \rho\le\frac14.
\end{equation}
The retained analytic finite-vertex estimate asserts that, for every
formal finite forest of every rank $n\ge59$,
\[
 \beta_n-\bigl(q+2\xi(u_{16}+u_{17})\bigr)\cdot\Phi_n
                                                       \ge\frac14.
\]
Since finite raw features are nonnegative and feature 2 is zero,
\eqref{fv:finite-tail-rows} proves its vertex budget. This is a
spectator- and rank-uniform analytic input, separate from the rank-58
catalogue. The coordinatewise comparison is a sufficient restriction
on the searched tail law, not a necessary description of all valid
finite-vertex charges.

\subsection{Optimization of the finite-vertex budget}
\label{fv:377}

The finite-vertex problem asks for the smallest allowance required
by a charge family. Once a finite formal forest domain $\mathcal T$ and a
local coefficient family $\mathcal C$ have been specified, write
\begin{equation}\label{fv:budget-template}
 U_{\mathcal T}(a)=\max_{F\in\mathcal T}a\cdot\Phi(F),
 \qquad \min_{a\in\mathcal C}U_{\mathcal T}(a).
\end{equation}
Here $a$ is the coefficient vector used at the finite vertex,
$\Phi(F)$ is its aggregate feature vector, and $\mathcal C$ is the
projection of the specified joint certificate family onto $a$.
The auxiliary variables and their conditions remain part of that
joint family.
The relevant geometric allowance is supplied independently by
incidence. A complete feasible certificate closes the finite-vertex
part when its required allowance is no greater than that geometric
allowance. Its face and ideal-vertex obligations, including all
dependencies of their local proofs, must hold at the same shared
coefficients.

Here $\mathcal T$ is the complete admitted rank-58 domain, and the
independently supplied allowance is $\beta_{58}=377$. The following
specialization gives both a feasible law and an exact dual lower bound
for its declared coefficient family. It makes no extension of the
endpoint catalogue to other ranks.

The following rational point uses the coordinate order specified in
Section~\ref{fv:features}:
\begin{align}
 (q_0,\ldots,q_{12})={}&
 (0,1,1,1,1,\tfrac12,\tfrac14,\tfrac14,\tfrac14,\tfrac14,
       \tfrac58,\tfrac58,\tfrac12),\notag\\
 (q_{13},\ldots,q_{25})={}&
 (\tfrac12,\tfrac12,\tfrac12,-\tfrac12,-\tfrac12,\tfrac12,
       \tfrac18,\tfrac14,\tfrac12,\tfrac38,\tfrac58,\tfrac12,
       \tfrac14),\notag\\
 a^{58}={}&q,\qquad
 a^{\rm tail}=q+\tfrac12(u_{16}+u_{17}),\notag\\
 \lambda={}&\tfrac12,\qquad\rho=0.
 \label{fv:successful-law}
\end{align}
In particular the two $D_{13}$ indicator coefficients are negative at
the endpoint and zero on the tail. This does not contradict the
nonnegativity of the resulting corner charge, which is checked on
the admitted decorated-corner domain.

The private variables can also be specified exactly. In the finite-tail
and descent blocks take $\xi=1/4$ and $\mu=13/2187$. Take
$k_r=\max_{\operatorname{rank}\Gamma=r}q_\Gamma$, over the complete
connected-family inventory, and use the retained ideal payoff vectors
and their ancestry equalities to assign every $B_{R,i,a}$, with
$\tau=0$. This is an assignment to the complete ideal certificate,
not a selection of a numerically best product profile. The actual
complete verification also checks every finite-network potential at
these same shared coefficients.

For example, the face rows
\eqref{fv:example-face-P}--\eqref{fv:example-face-M} take values
$17/16$ and $1$ on both branches. The guarded ideal rows have slack
at least $939/13312$ in type I and $50623/25792$ in type II. These
examples illustrate the assignment; complete exact verification checks
every local and network row at the same shared coefficients.

There is a useful optimization statement in addition to feasibility.
Let $\mathcal A$ be the projection onto
$(a^{58},a^{\rm tail},\lambda)$ of the complete local coefficient
system, including \eqref{fv:A13-rows}, with $\rho=0$. Replace the
last network row by $U_S\le U$, for one free common variable $U$,
and minimize $U$. Eliminating the potentials gives
\begin{equation}\label{fv:minimax}
 \min_{(a^{58},a^{\rm tail},\lambda)\in\mathcal A}
                 \max_F a^{58}\cdot\Phi(F),
\end{equation}
where $F$ ranges over the complete admitted rank-58 formal-forest
domain. The allowance 377 is no longer imposed as an upper bound
in this diagnostic optimization.

\begin{proposition}[Exact restricted minimax]\label{fv:exact-minimax}
The optimum in \eqref{fv:minimax} is 377.
\end{proposition}

\begin{proof}
The complete feasible assignment \eqref{fv:successful-law}, including
its private local variables and all network potentials, gives an upper
bound of 377. For the converse, write $a_0,a_1$ for coordinates 0 and
1 of $a^{58}$. The retained corner and triangle rows give
$a_0\ge0$ and $3a_0+2a_1\ge2$.
The formal forest
\[
                    F_*=2D_{17}+D_{24}
\]
has total rank 58. The ordinary table gives
$S_{D_{17}}=227/2$ and
$S_{D_{24}}=8\cdot24-32-1-4-3-2=150$. All ordinary cross
and contextual features vanish on this forest, so its aggregate
feature vector is
$\Phi(F_*)=1653u_0+377u_1$. Thus one path bound gives
$U\ge1653a_0+377a_1$. The exact nonnegative combination
\begin{equation}\label{fv:minimax-identity}
 U-377=(U-1653a_0-377a_1)
 +\frac{377}{2}(3a_0+2a_1-2)+\frac{2175}{2}a_0
\end{equation}
proves the lower bound: the first term is the sum of this forest's
network rows and its link to $U$, and the other terms are multiples
of the two local rows.
\end{proof}

The allowance 377 comes from incidence; the optimum 377 comes from
feasibility and \eqref{fv:minimax-identity}. This is optimality within
the declared law family and formal domain. The forest $F_*$ need not
be geometrically realized, and richer valid law families may have
smaller optima.

In the original capped system,
$\max_Fa^{58}\cdot\Phi(F)+\rho\le377$ and $\rho\ge0$.
The same inequalities force $a_0=0$, $a_1=1$ and $\rho=0$.
Thus maximizing this margin and requiring it to be positive would
miss the successful certificate. Incidence already has the strict
excess needed for the proof.

\subsection{All dimensions and the recovered bound}
\label{fv:all-dimensions}

For clarity, the separate coverage obligations are collected below.
They use one dimension-dependent charge, not independent face and
vertex laws.
\begin{center}\small
\begin{tabular}{p{.29\textwidth}p{.61\textwidth}}
\toprule
Domain&Retained certificate argument\\\midrule
Finite vertices, $n=58$&The complete vector-cost network and all its
 entry and transition inequalities.\\
Finite vertices, $n\ge59$&The analytic finite-vertex estimate and
 coordinatewise domination \eqref{fv:finite-tail-rows}.\\
Ideal vertices, $n=58$, $\chi_B$ aggregate zero&The complete rank-57
 ideal block and connected-family links.\\
Ideal vertices, $n=58$, positive $\chi_B$ aggregate&The complete
 deletion classification \eqref{fv:guarded-links} and all thirteen
 rows \eqref{fv:guarded-rows}.\\
Ideal vertices, $n=59$, largest factor rank 10--16&The same costs
 and the complete rank-descent inequalities.\\
Other ideal vertices at 59, and all $n\ge60$&The analytic ideal
 surplus theorem and the four rows \eqref{fv:analytic-ideal-rows}.\\
All actual corners and small faces, $n\ge58$&The complete local
 geometric and contact coverage of Input~\ref{fv:face-assumption}.\\
\bottomrule
\end{tabular}
\end{center}

\begin{theorem}[Finite-volume recovery]\label{fv:recovery}
The retained local geometric, classification, incidence and exact
finite-data theorems, together with the verified common-coefficient
certificate above, exclude every finite-volume hyperbolic Coxeter
polytope in every integer dimension $n\ge58$.
\end{theorem}

\begin{proof}
Suppose such a polytope exists. The last row of the table gives
nonnegative corner charges and the required payments on its triangles
and quadrilaterals. The preceding rows give
\[
 \sum_{\kappa\text{ at }v}w(\kappa)\le b_n(v)
\]
at every finite or ideal vertex. Apply the charge certificate and the
actual lift of Lemma~\ref{fv:actual-lift}.
\end{proof}


\section{Energy budgets and the common capacity dual}
\label{sec:energy-framework}

Both compact cases use the same specialization of
Theorem~\ref{thm:homogeneous-certificate}. A complete domain $b$ has
its own original orthogonal energy slots and all applicable rows
\begin{equation}\label{eq:branch-energy}
                  A_be_b\ge c_b,\qquad e_b\ge0.
\end{equation}
Here $A_b\ge0$ is rational, each row has a positive entry, and
$c_b\ge0$ is rational. In the actual matrices $A_b$ records supports
by $0$ and $1$. More generally, let the fixed geometric budget be
\begin{equation}\label{eq:general-energy-budget}
              h_b^Te_b<B_b,\qquad h_b>0,\quad B_b>0.
\end{equation}
The weights $h_b$ and budgets $B_b$ are rational data.
The compact geometry gives $h_b=\one$ and $B_b=1$ by
\eqref{eq:component-energies}. Leaving them as data shows which part
of the formulation is independent of that numerical specialization.

Different complete domains retain separate coordinates. For example,
a two-slot configuration and a three-slot configuration do not share
energy variables merely because both occur in the cover. Within one
domain, however, every use of a given original component refers to
its one slot. This distinction preserves both local applicability and
energy ownership.

\subsection{One normalization and duality theorem}

For each domain introduce nonnegative mass $\tau_b$, scaled energies
$E_b$ and deficit $\delta_b$. Impose
\begin{equation}\label{eq:budget-deficit}
 \begin{gathered}
 A_bE_b\ge c_b\tau_b,\qquad
 \delta_b=B_b\tau_b-h_b^TE_b\ge0,\\
                         \sum_b\delta_b=1.
 \end{gathered}
\end{equation}
A proved local exclusion adds
\begin{equation}\label{eq:excluded-mass}
                              \tau_b=0.
\end{equation}
It then forces $E_b=0$ and $\delta_b=0$, since $h_b>0$.
A fixed positive tolerance is unnecessary: the positive physical
budget deficit itself supplies the normalization.

\begin{theorem}[Budget lift and capacity dual]
\label{thm:budget-lift}\label{prop:general-energy-budget}
\label{thm:energy-normalization}\label{thm:energy-certificate}
\label{prop:capacity-complete}
For a finite complete domain family, the following assertions hold.
\begin{enumerate}[label=\textup{(\roman*)}]
\item System~\eqref{eq:budget-deficit}--\eqref{eq:excluded-mass} is
feasible exactly when a nonexcluded domain admits energies satisfying
\eqref{eq:branch-energy} and~\eqref{eq:general-energy-budget}.
Every actual polytope in the covered dimensions therefore has a lift.
\item Put
\begin{equation}\label{eq:branch-minimum}
 m_b=\min\{h_b^Te:e\ge0,\ A_be\ge c_b\}.
\end{equation}
Then
\begin{equation}\label{eq:branch-duality}
 m_b=\max\{c_b^T\alpha:\alpha\ge0,\ A_b^T\alpha\le h_b\},
\end{equation}
and both optima are attained rationally. The normalized system is
infeasible exactly when every
nonexcluded domain has $m_b\ge B_b$.
\item Consequently nonnegative rational multipliers satisfying
\begin{equation}\label{eq:general-capacity-certificate}
               A_b^T\alpha_b\le h_b,\qquad
               c_b^T\alpha_b\ge B_b
\end{equation}
for each nonexcluded domain give a complete infeasibility certificate
for this linear energy model.
\end{enumerate}
\end{theorem}
\begin{proof}
For an energy vector with $s=B_b-h_b^Te_b>0$, take
$\tau_b=1/s$, $E_b=e_b/s$, $\delta_b=1$ and set all other domains
zero. Conversely, some deficit in a normalized feasible point is
positive. Its mass is positive, so division by that mass gives the
local rows and
\[
           h_b^T(E_b/\tau_b)=B_b-\delta_b/\tau_b<B_b.
\]
Completeness of the geometric cover proves the actual-object claim.
The masses can be arbitrarily large; no artificial upper bound is
imposed on them.

For the optimum, a sufficiently large multiple of $\one$ satisfies
all local rows. Positivity of $h_b$ makes every finite objective
sublevel compact, so the minimum is attained. To see the dual signs
and equality directly, omit the subscript and consider the polyhedral
cone
\[
       K=\{(Ae-u,h^Te+t):e,u,t\ge0\}.
\]
It contains $(c,m)$ but no $(c,m-\varepsilon)$ for
$\varepsilon>0$. A nonzero supporting functional $(-\alpha,\beta)$,
nonnegative on $K$ and zero at $(c,m)$, gives
\[
 \alpha\ge0,\quad\beta\ge0,\quad A^T\alpha\le\beta h,
                  \quad\alpha^Tc=\beta m.
\]
If $\beta=0$, the positive entry in each nonnegative row of $A$
forces $\alpha=0$, a contradiction. Division by $\beta$ gives a
dual optimizer. The reverse inequality follows by multiplying
$Ae\ge c$ by any feasible dual vector. A vertex of the compact
primal optimal face is a rational vertex of the original feasible
polyhedron. The resulting dual optimal face is also a nonempty
rational polyhedron, so has a rational point by
Lemma~\ref{lem:cone-alternative}.

Finally, multiplication of the scaled rows by a capacity certificate
shows
\[
 h_b^TE_b\ge\alpha_b^TA_bE_b
            \ge(c_b^T\alpha_b)\tau_b\ge B_b\tau_b.
\]
Thus every deficit is zero. In the common formulation, take $C_b$ to
be the local cone in~\eqref{eq:budget-deficit}, with zero mass on
excluded domains, $T_b=0$ and $d_b=\delta_b$.
The displayed inequality is precisely the certificate
$-d_b\ge0$ of Theorem~\ref{thm:homogeneous-certificate}.
Part (ii) makes this certificate criterion necessary as well as
sufficient for the stated finite linear model.
\end{proof}

\subsection{Capacities and optimization in the compact applications}

For $h_b=\one$, $B_b=1$, the implemented dual seeks
\begin{equation}\label{eq:capacity-dual}
 \alpha_b\ge0,\qquad A_b^T\alpha_b\le\one,\qquad
              c_b^T\alpha_b\ge\eta,\qquad\eta\ge1.
\end{equation}
There is one capacity inequality for each original energy slot. For
support rows its left side is
$\sum_{\ell:j\in S_\ell}\alpha_{b\ell}$: the total multiplicity with
which that slot is charged. A new local theorem does not give another
copy of the same energy.

\begin{example*}
Suppose $e_I+e_A\ge a$ and $e_I+e_B\ge b$. Weights one on both rows
would count $e_I$ twice. Weights $1/2,1/2$ give instead
\[
           e_I+\tfrac12e_A+\tfrac12e_B\ge\tfrac12(a+b),
\]
a valid lower bound for the total energy because the omitted
half-energies are nonnegative. Capacities test validity of the
combination; its payment tests whether it reaches the budget.
\end{example*}

For one fixed dimension-labelled catalogue define
\begin{equation}\label{eq:threshold-energy-value}
 \Gamma_N=\min_{\substack{b\ \text{not locally excluded}\\ n_b\ge N}}m_b,
 \qquad\min\varnothing=+\infty.
\end{equation}
The normalized LP is infeasible exactly when $\Gamma_N\ge1$.
Raising $N$ removes domains, so this minimum cannot decrease; this
is the energy form of Proposition~\ref{prop:threshold}.
Completeness concerns the supplied linear inequalities, not geometric
realizability of an arbitrary feasible energy vector.

The recorded computation enforces~\eqref{eq:capacity-dual} and
maximizes the sum of local payments, favoring simple saturated
capacity solutions. Exact verification of the resulting rational
point proves the needed comparison. No numerical optimality claim
is used. Compactness and compact arithmeticity change the domains
and rows; their linear assembly is identical.

\section{The compact bound from a uniformly generated family}
\label{c:section}

The compact instance uses the component theorems of Supplement~C
of~\cite{BoundsSupplement}. We state their energy bounds and complete
domains; the common LP then performs the global accounting.

Fix the original farthest vertex and the original vector $x$ used in
Theorem~\ref{thm:energy-normalization}. If $C$ is an irreducible spherical
component of its vertex diagram, put
\[
 e_C=\norm{\operatorname{proj}_{\langle C\rangle}x}^2.
\]
The component subspaces are mutually orthogonal and exhaust the spatial
space. Consequently
\begin{equation}\label{c:budget}
 e_C\ge0,\qquad \sum_C e_C=\norm{x}^2<1.
\end{equation}
For disjoint collections $\mathcal S$ of components we write
$e_{\mathcal S}=\sum_{C\in\mathcal S}e_C$. Throughout this section, a
\emph{spectator} is another component at this same vertex, with its actual
opposite walls retained. It is not an independently chosen diagram.

Vinberg's compact dimension theorem~\cite{Vinberg1984} excludes dimensions
at least $30$. The LP covers each actual dimension $18\le n\le29$;
no passage to a lower-dimensional face is used.

\subsection{Small components and the actual exceptional profiles}
\label{c:small}

We use the usual notation for irreducible spherical Coxeter diagrams.
The rank-two family is $I_2(m)$, where $m\ge3$ is an integer. We identify
$A_2$ with $I_2(3)$ and do not count aliases such as $B_2$ or $D_3$ a
second time. An \emph{ordinary} component is a component that does not
carry one of the actual exceptional wall profiles described below.

A profile specifies an opposite wall of the given polytope and its
contacts with the original walls. An old bond labelled seven alone
does not specify $O_7$. All labels below belong to pairs of actual
original walls. Auxiliary projected vectors are not assigned those labels.

\begin{proposition}[Ordinary component estimates, retained input]
\label{c:ordinary-input}
For the same original vector $x$, the ordinary components in
Table~\ref{c:ordinary-table} have energy strictly greater than the listed
number. Every estimate permits arbitrary finite spectators.
\end{proposition}

\begin{table}[htbp]
\centering
\caption{The ordinary component lower bounds. Each decimal in this paper
is an exact rational number.}\label{c:ordinary-table}
\begin{tabular}{ll@{\qquad}ll}
\toprule
Component&Lower bound&Component&Lower bound\\\midrule
$A_1$&$1/10$&$A_2$&$319/2000$\\
$I_2(m),\ m\ge4$&$1/7$&$A_3$&$1/4$\\
$B_3,H_3$&$2/7$&$A_4$&$2/7$\\
$B_4,D_4,F_4$&$1/3$&$H_4$&$11/36$\\
$A_5$&$3/8$&$D_5$&$5/14$\\
$A_6,B_5,B_6,D_6$&$3/7$&$B_7,D_7$&$1/2$\\
$B_8,D_8$&$4/7$&$E_6$&$9/20$\\
$E_7$&$5/9$&$E_8$&$2/3$\\\bottomrule
\end{tabular}
\end{table}

These are the ordinary-component conclusions of Supplement~C, including
the $B_3,H_3$ improvement. The dihedral row quantifies over every integer
$m\ge4$, rather than a finite sample of labels.

For clarity, we describe the exceptional profiles before using their
energy debits. Let $p$ be the original unit timelike vertex vector, and
let $u$ be the unit normal of an original $A_1$ component. An isolated
short rank-one opposite has the form
\begin{equation}\label{c:I-normal}
 c_I=\sqrt h\,p-\sqrt{1+h}\,u,
 \qquad 0<h<\frac12.
\end{equation}
Its spatial support is this one $A_1$. This is profile $I$.
Profile $J$ is the distinct short-wall alternative supported on two old
$A_1$ factors, with a dotted contact on one and an actual contact labelled
three on the other. Its squared time lies in $(1/3,3/8)$. The two factors
belong to one profile and must be accounted for jointly.

There are also two rank-two profiles:
\[
 N_7=(\text{old }A_2,\ \text{actual new isolated contact }7),
 \qquad
 O_7=(\text{old }I_2(7),\ \text{actual new isolated contact }3).
\]
Here isolation means that the new wall has no spatial contact with
another old component. These specifications include the new wall, not
merely the type of its old support.

\begin{proposition}[Exceptional profiles and compatibility, retained input]
\label{c:defect-input}
At most one of $I,J,N_7,O_7$ occurs, and the alternatives do not coexist.
Their owned ranks and energy bounds are
\[
\begin{array}{c|c|c}
\text{profile}&\text{old support}&\text{energy lower bound}\\\hline
I&A_1&0\\
J&A_1+A_1&17/50\\
N_7&A_2&9/100\\
O_7&I_2(7)&9/100.
\end{array}
\]
The positive entries are strict in the geometric theorem. No owned
component receives an ordinary bound in addition to its profile bound.
\end{proposition}

Compatibility is part of the actual-wall theorem, including overlapping
short-wall alternatives; it is not inferred from the spherical types
of the supports.

Two further input families will be used. If $T_j$ is an ordinary
irreducible of rank $j$, then
\begin{align}
 e_{I+A_1}&>1/5,&
 e_{I+2A_1}&>3/10,\label{c:I-joint-one}\\
 e_{I+T_2}&>21/100,&
 e_{I+T_3}&>2/7.\label{c:I-joint-two}
\end{align}
Each left side is one joint support. Overlapping uses of $I$ therefore
share its single capacity in Theorem~\ref{thm:energy-certificate}.

Beside an actual $N_7$ or $O_7$, every other ordinary rank-two component
has energy greater than $1/5$. An additional ordinary $A_1$ has energy
greater than $1/5$ beside $N_7$, and greater than $2/15$ beside $O_7$.
These propagated estimates concern the additional component alone.
Their domains and the joint bounds above are those of the
ordinary-component appendix of Supplement~C.

\subsection{The larger component families}
\label{c:large}

The $A$-component theorems give
\begin{equation}\label{c:A-thresholds}
 e_{A_r}>\frac12\quad(r\ge7),\qquad
 e_{A_r}>\frac47\quad(r\ge8),\qquad
 e_{A_r}>\frac57\quad(r\ge13).
\end{equation}
All three estimates allow arbitrary finite spectators. The exact-rank
improvements retained by the LP are listed in
Table~\ref{c:A-table}. These are bounds on the original projection,
even when their local proofs use points or neighboring vertices.

\begin{table}[htbp]
\centering
\caption{Additional spectator-uniform $A$-component bounds.}
\label{c:A-table}
\begin{tabular}{clp{.57\linewidth}}
\toprule
Rank&Lower bound&Local result in Supplement~C\\\midrule
12&$13/20$&The two-end $A_{12}$ estimate\\
15&$413/500$&The $A_{15}$ reference-vector estimate\\
16&$43/50$&The $N=17$ reference-vector estimate\\
17&$6/7$&The $A$-chain central-window argument\\
18&$91/100$&The $N=19$ chain window and reference-vector estimate\\
19&$563/625$&The $A_{19}$ projection-energy lemma\\
20&$.92274$&The $A_{20}$ projection-energy lemma\\\bottomrule
\end{tabular}
\end{table}

For actual $B_r,D_r$ components with $r\ge9$, the common bound is
\begin{equation}\label{c:BD-nine}
                         e_{B_r},e_{D_r}>\frac9{14}.
\end{equation}
The next family uses an initial segment of the original component. Its
rank conditions are part of the statement and will be retained literally.

\begin{proposition}[The generic prefix estimate, retained input]
\label{c:prefix-input}
Let $6\le M\le16$. Suppose that the actual component is $B_r$ with
$r\ge M+5$, or $D_r$ with $r\ge M+6$. Put
\[
 f_i=\sqrt i-\frac{61}{100\sqrt2},\quad
 E_M=\frac{f_5^2}{5}+
        \sum_{i=5}^{M-1}(\sqrt{i+1}-\sqrt i)^2,
 \quad H_M=\left\lfloor\frac{M-5}{2}\right\rfloor.
\]
On the same original vector, the first $M+1$ coordinates of that
component have energy greater than
\begin{equation}\label{c:prefix-formula}
                  E_M-H_M/3200+3969/800000.
\end{equation}
The estimate permits arbitrary finite spectators.
\end{proposition}

The first two terms are the generic prefix estimate. The last term is
the additional initial-coordinate gain proved under exactly the same
hypotheses. It replaces the weaker bound on that prefix; it is not an
independent energy contribution. The rational lower bounds used in the
current program are $P_M+3969/800000$, with $P_M$ as follows:
\begin{center}
\begin{tabular}{c|ccccc}
$M$&7&8&9&10&11\\\hline
$P_M$&$.7349975$&$.7683275$&$.797425$&$.823735$&$.8472225$\\[2pt]
\end{tabular}
\par\smallskip
\begin{tabular}{c|ccccc}
$M$&12&13&14&15&16\\\hline
$P_M$&$.8689525$&$.88864$&$.90715$&$.9240775$&$.9401975$
\end{tabular}
\end{center}
The inequalities $P_M<E_M-H_M/3200$ are the rational evaluations in
Supplement~C. The current finite LP uses $7\le M\le16$ from this family.
For every applicable $M$ it may keep both the weaker and the stronger
row. On identical support and domain the weaker one is redundant, but
retaining it does not alter the validity of the stronger row.

There are three further local continuations. They are stated separately
because their proofs extend a short end where the generic theorem alone
does not apply:
\begin{enumerate}[label=\textup{(\roman*)}]
\item On the entire forest $B_{16}+T_2$ or $D_{16}+T_2$ in actual
 dimension $18$, where $T_2$ is ordinary, the large component has energy
 greater than $.8689525$.
\item On the entire forest $B_{17}+A_1$ in dimension $18$, with ordinary
 nonshort $A_1$, the $B_{17}$ component has energy greater than $.90715$.
\item On the entire forest $D_{19}+A_1$ in dimension $20$, with ordinary
 nonshort $A_1$, the actual $D_4$ continuation and initial gain give
 $e_{D_{19}}>.91211125$.
\end{enumerate}
The first two are the exact-rank-eighteen layer estimates. The third
uses the proper $D_{19}$ continuation before the rank-twenty accounting.
We import that local continuation and its energy calculation, not the
conclusion that dimension twenty is impossible. A bound proved for an
ordinary $T_2$ is not used for $N_7$ or $O_7$, and an ordinary $A_1$
bound is not transferred to $I$.

\subsection{A complete domain construction}
\label{c:cover}

We call $A_r$, $r\ge7$, and $B_r,D_r$, $r\ge9$,
\emph{dominant}. This is a definition for the present accounting, not
a new class of Coxeter diagrams. Every dominant costs more than one
half by \eqref{c:A-thresholds} and \eqref{c:BD-nine}. Every other
ordinary irreducible costs at least its rank divided by fourteen, as
one checks directly from Table~\ref{c:ordinary-table}.

\begin{lemma}[The common spectator estimate]\label{c:spectator}
Let a collection of nondominant components have total rank $S$. Include
every exceptional profile owned by that collection. Its energy satisfies
\begin{equation}\label{c:spectator-rate}
                          e_{\mathcal S}\ge\frac{S-1}{14}.
\end{equation}
For no profile or profile $J$, the stronger bound $S/14$ is available.
For profile $I$ the possible loss is $1/14$, and for $N_7$ or $O_7$
it is $37/700$.
\end{lemma}
\begin{proof}
Sum the ordinary rank-over-fourteen estimates on the unmarked components.
Profile $I$ replaces the ordinary rank-one contribution by a nonnegative
energy, losing at most $1/14$. Profile $J$ owns rank two and its bound
$17/50$ is greater than $2/14$, so it creates no loss. Either typed
rank-two profile has cost at least $9/100$, whose deficit from $2/14$ is
$37/700<1/14$. Proposition~\ref{c:defect-input} permits at most one of
these losses. This proves every assertion.
\end{proof}

For a fixed dimension $18\le n\le29$, the following construction produces
a complete necessary family of energy domains.

\begin{enumerate}[label=\textup{\arabic*.},leftmargin=*]
\item If there is no dominant, use one aggregate energy variable for all
 components and the row $e\ge(n-1)/14$.
\item If there are at least two dominants, select two distinct actual
 dominant components. Give each its own energy variable and the bound
 $1/2$. A third aggregate variable holds every remaining component.
 The selected components' individual ranks need not be fixed for this
 domain: only their disjointness and their common lower bound are used.
\item If there is exactly one dominant of type $T$ and rank $r$, let
 $S=n-r$ be the total spectator rank. Retain an explicit local
 component exclusion whenever one applies. Otherwise, for $S\ge8$ use
 the two rows $e_T\ge1/2$ and $e_{\mathcal S}\ge(S-1)/14$.
\item For $0<S<8$, let $d(T,r)$ be the largest applicable lower bound
 among the spectator-uniform component estimates in
 Section~\ref{c:large}. Only unconditional component estimates are used
 to choose $d(T,r)$; an estimate requiring one particular spectator
 does not qualify. If
 \begin{equation}\label{c:aggregate-test}
                  d(T,r)+(S-1)/14\ge1,
 \end{equation}
 retain the complete domain with the two aggregate rows giving these
 two bounds.
\item In every remaining case, retain every nondominant spherical
 spectator forest of rank $S$ and every allowed assignment of one
 exceptional profile. On that detailed domain include all applicable
 ordinary, component, joint and propagated bounds, together with the
 precisely matched local exclusions described in Section~\ref{c:local-exits}.
 For $S=0$ this step simply gives the whole dominant component.
\end{enumerate}

Coarse domains remain in the LP. Their displayed rows have a certificate
with weights one, including the equality case $S=8$. The test
\eqref{c:aggregate-test} decides whether to retain detailed spectator
types; it does not enlarge any local theorem's domain.

\begin{proposition}[Completeness of the compact family]
\label{c:complete-cover}
Every hypothetical compact Coxeter polytope of actual dimension
$18\le n\le29$, with its chosen original farthest vertex, gives either
one of the explicitly retained local contradictions or an energy point
in a domain of the above construction. Every coordinate is an energy
of the original vector, and every row applies on the entire domain to
which it is attached.
\end{proposition}
\begin{proof}
Spherical classification gives zero, at least two, or exactly one
dominant component. The first two alternatives give steps one and two.
In the third, its type and rank fix $S$; Lemma~\ref{c:spectator} and the
applicable component bound justify either coarse description.
Otherwise step five includes the actual spectator multiset and its
unique allowed profile. Equal-type components may be relabelled with
their actual walls to put a marked component first. This preserves all
other parameters: an ordinary dihedral component retains its full
integer range, while an $O_7$ mark fixes only its owned component to
label seven. Every row and local exclusion is then used with its stated
component, profile and spectator hypotheses. Orthogonality gives
\eqref{c:budget} on the original vector.
\end{proof}

To generate a detailed remainder, use this ordered list of nondominant
types through spectator rank seven:
\[
\begin{array}{c|l}
\text{rank}&\text{types}\\\hline
1&A_1\\
2&A_2=I_2(3),\ I_2(m)\ (m\ge4)\\
3&A_3,B_3,H_3\\
4&A_4,B_4,D_4,F_4,H_4\\
5&A_5,B_5,D_5\\
6&A_6,B_6,D_6,E_6\\
7&B_7,D_7,E_7.
\end{array}
\]
Choose successive types in nondecreasing order, subtracting their ranks
until the remainder is zero. Then mark none, one $A_1$ as $I$, two
$A_1$ factors as $J$, one $A_2$ as $N_7$, or one dihedral factor as
$O_7$. Nondecreasing order and exchangeability remove repetitions,
while all dihedral parameters remain quantified. The construction
produces the $829$ domains checked below; its completeness follows
from Proposition~\ref{c:complete-cover}, independently of that count.

\subsection{The retained local contradiction domains}
\label{c:local-exits}

Table~\ref{c:local-exit-table} gives the local contradiction inputs of
Supplement~C. Their normalized mass is zero; no component energy floor
is inferred from them. ``Entire'' specifies the complete original
vertex forest. Apart from C1, only the displayed entire-forest domains
are used, with every Coxeter label and separated-wall parameter in the
cited theorem retained.

\begin{longtable}{>{\raggedright\arraybackslash}p{.08\linewidth}
 >{\raggedright\arraybackslash}p{.43\linewidth}
 >{\raggedright\arraybackslash}p{.39\linewidth}}
\caption{Compact local contradiction inputs. The descriptions refer to
the complete local proofs in Supplement~C.}\label{c:local-exit-table}\\
\toprule
Code&Matched domain&Local proof retained\\\midrule
\endfirsthead
\toprule Code&Matched domain&Local proof retained\\\midrule\endhead
C1&Any $A_r$, $21\le r\le29$; $B_r$, $23\le r\le29$;
 or $D_r$, $24\le r\le29$, with arbitrary finite spectators&
 Large-component cap theorem\\
C2&Entire $B_r+I$, $r\ge20$, or $D_r+I$, $r\ge21$&
 The uniform proper classical-component argument beside $I$\\
C3&Entire $A_{16}+N_7$ or $A_{16}+O_7$&
 The two typed endpoint families in actual dimension $18$\\
C4&Entire $A_{17}+A_1$ or $A_{17}+I$, with ordinary $A_1$&
 The complete rank-one-remainder argument in dimension $18$\\
C5&Entire $B_{16}+N_7/O_7$ or $D_{16}+N_7/O_7$&
 The typed rank-sixteen arguments, retaining both signed $D$ forks\\
C6&Entire $B_{17}+I$&The short rank-one-remainder head argument\\
C7&Entire $D_{17}+A_1$ or $D_{17}+I$, with ordinary $A_1$&
 The separate ordinary and isolated-short rank-one arguments\\
C8&Entire $A_{18}$&The two-end argument using the gap between
 allowed actual wall products\\
C9&Entire $B_{18}$&The full terminal-wall partition, including common
 and distinct opposite walls\\
C10&Entire $D_{18}$&The complete original end-wall and fork argument,
 including all finite, parallel and dotted contacts at the distinguished end\\
C11&Entire $A_{17}+N_7/O_7$ or $A_{18}+I$&
 The separate exact-rank-nineteen local arguments\\
C12&Entire $D_{17}+T_2$, $B_{17}+N_7/O_7$, or
 $B_{18}/D_{18}+A_1$, with ordinary $T_2,A_1$&
 The proper classical rank-nineteen continuation\\
C13&Entire $D_{17}+N_7/O_7$, $B_{18}+I$, or $D_{18}+I$&
 The respective typed and short-remainder arguments at rank nineteen\\
C14&Entire $A_{19}$, $B_{19}$ or $D_{19}$&
 Their separate complete endpoint, head and fork proofs\\
C15&Entire $A_{19}+I$&The $A_{19}$ isolated-short-remainder argument\\
C16&Entire $B_{18}/D_{18}+N_7/O_7$ or $B_{19}/D_{19}+I$&
 The proper classical rank-twenty continuation\\
C17&Entire $A_{20}$, $B_{20}$ or $D_{20}$&
 Their separate whole-component arguments at actual rank twenty\\
C18&Entire $A_{19}+N_7/O_7$ or $A_{20}+I$&
 The typed and isolated-short $A$-component arguments\\
C19&Entire $D_{19}+N_7/O_7$ or $D_{20}+I$&
 The proper $D$-component continuation at rank twenty-one\\
C20&Entire $B_{21}$ or $D_{21}$&
 The separate full-component arguments at actual rank twenty-one\\
C21&Entire $B_{22}$, $D_{22}$ or $D_{23}$&
 The retained full-component terminal argument\\\bottomrule
\end{longtable}

Each entry denotes the complete cited local theorem: C10, for example,
includes all whole-$D_{18}$ end-wall and fork alternatives. A comparison
with another vertex within such a proof contributes no further energy
to the original budget \eqref{c:budget}.

\subsection{The resulting compact certificate}
\label{c:certificate}

The construction gives $829$ branch domains. Of these, $151$ carry one
of the retained local contradictions. On each remaining branch the LP
has one nonnegative multiplier for every applicable energy row. The
capacity at each original projection is one. The exact rational
certificate satisfies the hypotheses of
Theorem~\ref{thm:energy-certificate} simultaneously in every branch.
In particular, overlapping rows involving $I$ share the same $I$ column;
the program cannot buy two copies of its energy.

\begin{theorem}[Compact recovery]\label{c:recovery}
The local geometric inputs stated in Sections~\ref{c:small},
\ref{c:large} and~\ref{c:local-exits}, together with spherical
classification, give an infeasible normalized energy LP for the complete
union of actual dimensions $18,\ldots,29$. With Vinberg's external
compact ceiling $29$, they recover the dimension bound $n\le17$.
\end{theorem}
\begin{proof}
Proposition~\ref{c:complete-cover} assigns every hypothetical polytope
to a complete branch. A branch with a retained local contradiction is
impossible. On every other branch, the exact capacity certificate gives
total original energy at least one by
Theorem~\ref{thm:energy-certificate}, contrary to \eqref{c:budget}.
The exact check covers all $3157$ dual variables and $2224$ rows;
the reproduction section identifies the certificate and model records.
The external theorem excludes dimensions $n\ge30$.
\end{proof}

Thus the LP performs the final rank and energy accounting from valid
local bounds. Its domain construction begins with spherical
classification, without assuming the historical eleven-forest
rank-eighteen list.

\section{The compact arithmetic bound with the same energy program}
\label{a:section}

For a compact arithmetic polytope, its reflection group is an arithmetic
lattice. Supplement~A of~\cite{BoundsSupplement} supplies the further
local restrictions below, with the actual defining field and original
root lines retained. The energy budget remains \eqref{c:budget}.
The compact catalogue already covers dimensions $18,\ldots,29$;
we append the arithmetic domains at actual ranks $16$ and $17$ to
that same catalogue.

\subsection{The retained arithmetic inputs}
\label{a:inputs}

We first recall what the field conditions refer to. In the quadratic-form
model supplied by the arithmetic companion, $K$ is a totally real number
field, $q$ is a quadratic form over $K$, and $B$ is its associated symmetric
bilinear form. One real embedding $\sigma_0:K\hookrightarrow\R$ is
distinguished: at this embedding $q$ has Lorentzian signature and defines
the given hyperbolic space. At every other real embedding $\sigma$, the
conjugate form $q^\sigma$ is positive definite. This is a statement about
the form, not about independence of an arbitrary collection of walls.

The roots $\alpha_i$ span the original wall lines over this same field.
Their Cartan entries are
\begin{equation}\label{a:cartan}
 C_{ij}=\frac{2B(\alpha_i,\alpha_j)}{B(\alpha_i,\alpha_i)}.
\end{equation}
The arithmetic integrality condition used in the local proofs requires
the cyclic products
\begin{equation}\label{a:cycles}
             C_{i_1i_2}C_{i_2i_3}\cdots C_{i_\ell i_1}
\end{equation}
to lie in $\mathcal O_K$, the ring of algebraic integers in $K$.
This formulation is insensitive to changing a root representative on
its line: under $\alpha_i\mapsto t_i\alpha_i$, with $t_i\in K^\times$,
the entry $C_{ij}$ is
multiplied by $t_j/t_i$, and those factors cancel around the cycle.
The companion proves the field reductions on their stated domains.

At a nonphysical embedding $\sigma\ne\sigma_0$, a Gram matrix of
linearly independent conjugate roots is positive definite. A Gram matrix
of all wall roots can have linear dependencies and is only positive
semidefinite. In particular, one must not demand positive definiteness
of a full redundant wall matrix. Similarly, normalizing roots to unit
length can introduce square roots and enlarge the field. The field of
the resulting real cosine entries is not a substitute for the original
common field, root lines and norm classes. The local field-square tests
retain that information when they restrict the possible wall parameters.
Here a root's norm class is the square class of
$B(\alpha_i,\alpha_i)$ in $K^\times/(K^\times)^2$; rescaling by
$t_i\in K^\times$ changes its norm by $t_i^2$ and preserves that class.

The golden-lattice applications also retain their proved inclusion
$\Gamma(P)\subseteq\operatorname{Ref}(L)$ for the literal invariant
lattice $L$ of Bugaenko's theorem~\cite{Bugaenko1984}. Here $\Gamma(P)$
is the original reflection group and $\operatorname{Ref}(L)$ the full
reflection group of that specified lattice. Equality of defining
fields alone would not give this inclusion.

Unlike the compact construction above, the arithmetic instance begins
with proved necessary covers at ranks $16$ and $17$. Their pruning is
an input; the LP supplies their final energy accounting.

\begin{assumption}[Arithmetic input package]
\label{a:input-package}
We use the following proved statements of Supplement~A on their complete
stated domains:
\begin{enumerate}[label=\textup{(\roman*)},leftmargin=*]
\item The actual common totally real field, original root lines,
 integral Cartan-cycle conditions, field-square restrictions, and
 positivity at every nonphysical real embedding are retained wherever
 they occur in a local theorem.
\item The complete necessary rank-seventeen and rank-sixteen covers
 are the fourteen and nineteen groups stated in
 Tables~\ref{a:rank17} and~\ref{a:rank16}, including all defect,
 spectator and field alternatives used to prove those covers.
\item The component and joint-component energy theorems give the rows
 in Section~\ref{a:costs} and Tables~\ref{a:rank17}--\ref{a:rank16}
 on the specified supports and complete forest domains. Each component
 floor already includes all contributions used in its local proof.
\item The complete local contradictions in Table~\ref{a:local-exits}
 hold on their indicated actual domains, including their boundary
 cases and the identity of each actual rank-one opposite.
\item The four computed outputs in
 Table~\ref{a:computed-inputs} hold at their proved complete input
 domains, with the original reduction and domain-inclusion theorems.
\item Every golden-lattice application retains the specified lattice
 and reflection-group inclusion above.
\end{enumerate}
\end{assumption}

This package specifies proved inputs of Supplement~A, rather than
consequences of real matrix positivity alone.

\begin{table}[htbp]
\centering
\caption{The four retained computed-output premises. A surviving model
is a necessary algebraic possibility, not an asserted polytope.}
\label{a:computed-inputs}
\begin{tabular}{>{\raggedright\arraybackslash}p{.08\linewidth}
 >{\raggedright\arraybackslash}p{.43\linewidth}
 >{\raggedright\arraybackslash}p{.39\linewidth}}
\toprule
Output&Complete reduced domain&Conclusion used by the local proof\\\midrule
(i)&The quadratic-field whole-$A_{17}$ domain&Empty\\
(ii)&The two cubic-field whole-$A_{17}$ finite-anchor domains&Empty
 in each of the two complete domains\\
(iii)&The shared quadratic pair domain reached separately from whole
 $B_{17}$ and whole $D_{17}$&Two necessary pairs, each retained until
 the final $B$- or $D$-specific energy estimate is applied\\
(iv)&The quadratic pair domain reached from whole $B_{16}$&One necessary
 pair, retained until its actual-edge energy estimate is applied\\\bottomrule
\end{tabular}
\end{table}

The whole-$D_{16}$ exclusion is a separate manual dyadic argument,
not a census premise. Its domain also differs from $D_{16}+I$ at
rank seventeen; both appear separately below.

\subsection{Notation and the arithmetic local rows}
\label{a:costs}

All component energies in this section are still projections of the
original vector $x$. Write $T_2$ for an ordinary irreducible $I_2(m)$,
$m\ge3$, $T_3$ for $A_3,B_3$ or $H_3$, and $T_4$ for one of
$A_4,B_4,D_4,F_4,H_4$. Let $\mathcal K$ denote either actual typed profile
$N_7$ or $O_7$, with the two choices retained as distinct branches.
An unqualified $A_1$ is ordinary and nonshort. The symbols $I,J,N_7,O_7$
retain the actual-wall definitions of Section~\ref{c:small}.

We use the following complete rank-five families:
\begin{align*}
\mathcal R_5={}&\{A_5,D_5,A_4+A_1,H_4+A_1,
                   T_3+T_2,2T_2+A_1\},\\
\mathcal R'_5={}&\{A_5,D_5,A_4+A_1,H_4+A_1,
                   A_3+T_2,2T_2+A_1\}.
\end{align*}
Every member has ordinary energy at least $5/14$ by
Table~\ref{c:ordinary-table}. The restriction to $A_3$ in
$\mathcal R'_5$ belongs to the rank-sixteen input cover.

In addition to the ordinary estimates in
Table~\ref{c:ordinary-table}, the following joint bounds are used:
\begin{align}
 e_{I+T_2}&>21/100,& e_{I+T_3}&>2/7,&
 e_{I+T_4}&>23/70,\label{a:I-joints}\\
 e_{I+A_1}&>5051/25000=.20204.\label{a:I-A1}
\end{align}
These are the actual-short-profile theorems of Supplement~A.
We also retain the propagated floors $e_{T_2}>1/5$ beside either
typed profile and $e_{A_1}>2/15$ beside $O_7$.

Tables~\ref{a:rank17} and~\ref{a:rank16} state the remaining
arithmetic floors directly on their domains. Each entry is a retained
local theorem of Supplement~A for the complete forest in the middle
column, with every displayed dihedral and field alternative retained.
Separate summands denote separate support rows; a joint support is
named explicitly. A floor on a named large component concerns that
component alone, even when its local proof uses the remaining walls.
Thus $16/19$ on $A_{14}$ beside ordinary $T_2$ and
$.91286445198$ beside a typed profile are distinct input statements.
The ordinary $T_2$ floor $3/19$ covers $m\ge4$ by its arithmetic
dihedral theorem and $m=3$ by $319/2000>3/19$.
For typed $B/D$ components the single floors $.91859$ and $.93859$
already include the end-wall correction $.0487$; it consumes no
additional energy variable. Only the rows needed by the displayed
certificate are specified here; additional valid rows in the retained
model may receive weight zero.

\subsection{The complete rank-seventeen cover and its rows}
\label{a:rank17-section}

Table~\ref{a:rank17} gives all fourteen input groups. Braces and
semicolons denote unions of complete alternatives. X1--X13 are the
local contradiction inputs in Table~\ref{a:local-exits}.

\begin{longtable}{>{\raggedright\arraybackslash}p{.06\linewidth}
 >{\raggedright\arraybackslash}p{.35\linewidth}
 >{\raggedright\arraybackslash}p{.49\linewidth}}
\caption{All rank-seventeen groups and their LP rows.}
\label{a:rank17}\\
\toprule
No.&Complete forest family&Rows giving total energy at least one,
 or retained local contradiction\\\midrule
\endfirsthead
\toprule No.&Complete forest family&Energy rows or contradiction\\\midrule\endhead
1&$A_{11}+I+(A_5\text{ or }D_5)$&
 $13/20+5/14=141/140$; the nonnegative $I$ energy is unused\\
2&$A_{11}+\mathcal K+U_4$, $U_4\in\{A_4,B_4,D_4,F_4\}$&
 $13/20+9/100+2/7$\\
3&$A_{12}+R$, $R\in\mathcal R_5$&
 $5/7+5/14$\\
4&$A_{12}$ plus $I+T_4$; $I+T_3+A_1$; or $I+2T_2$&
 $5/7$ plus, respectively, the joint $23/70$;
 the joint $2/7$ and $1/10$; or the joint $21/100$ and $1/7$\\
5&$A_{14}$ plus $T_3$; $T_2+A_1$; or $I+T_2$&
 $4/5$ plus, respectively, $1/4$; $1/7+1/10$;
 or the joint $21/100$\\
6&$A_{15}$ plus $T_2$; $2A_1$; $I+A_1$; $N_7$; or $O_7$&
 $T_2$: X1. The next two give $.8348+.2$ and
 $.8348+.20204$. Either typed branch gives $.9191408+.09$\\
7&$A_{16}$ plus ordinary $A_1$ or actual $I$&X2 in each
 of its complete rank-one alternatives\\
8&Entire $A_{17}$&X3, retaining outputs (i) and (ii) and their
 actual endpoint-domain inclusions\\
9&$B_{15}+N_7$ or $B_{15}+O_7$&
 $.93859+.09=1.02859$\\
10&$B_{16}+I$&X4\\
11&Entire $B_{17}$&$e_{B_{17}}\ge1.00184$, after the
 separate reduction and energy estimates for both survivors\\
12&$D_{15}+N_7$ or $D_{15}+O_7$&
 $.91859+.09=1.00859$\\
13&$D_{16}$ plus ordinary $A_1$ or actual $I$&
 X5 for ordinary $A_1$; X6 for actual $I$\\
14&Entire $D_{17}$&$e_{D_{17}}\ge1.00184$, after its
 genuine-fork reduction and its energy estimates for both survivors\\\bottomrule
\end{longtable}

The fourteen groups yield $39$ complete parameter-family branches,
without asserting their geometric realizability. Repeated symbols
such as $2T_2$ represent independently parameterized components.
In $I+2T_2$, choose one dihedral for the joint bound and give the
other its ordinary row; the $I$ energy is used once.

\subsection{The complete rank-sixteen cover and its rows}
\label{a:rank16-section}

Table~\ref{a:rank16} gives the separate nineteen-group rank-sixteen
input cover, including its narrower $\mathcal R'_5$ family.

\begin{longtable}{>{\raggedright\arraybackslash}p{.06\linewidth}
 >{\raggedright\arraybackslash}p{.35\linewidth}
 >{\raggedright\arraybackslash}p{.49\linewidth}}
\caption{All rank-sixteen groups and their LP rows.}
\label{a:rank16}\\
\toprule
No.&Complete forest family&Rows giving total energy at least one,
 or retained local contradiction\\\midrule
\endfirsthead
\toprule No.&Complete forest family&Energy rows or contradiction\\\midrule\endhead
1&$A_{10}+I+(A_5\text{ or }D_5)$&$404/625+5/14$\\
2&$A_{10}+\mathcal K+U_4$, $U_4\in\{A_4,B_4,D_4,F_4\}$&
 $404/625+9/100+2/7$\\
3&$A_{11}+R$, $R\in\mathcal R'_5$&$13/20+5/14=141/140$\\
4&$A_{11}+I+T_4$&The joint $A_{11}+I$ bound $18/25$,
 plus the disjoint $T_4$ bound $2/7$, giving $176/175$\\
5&$A_{11}+O_7+A_3$&$13/20+9/100+13/50=1$\\
6&$A_{12}$ plus $A_4$; $H_4$; or $2T_2$&$5/7+2/7=1$\\
7&$A_{12}+I+T_3$&$5/7$ plus the joint $I+T_3$ bound $2/7$\\
8&$A_{12}+\mathcal K+T_2$&$5/7+9/100+1/5$, with the actual
 typed-profile propagation on the last component\\
9&$A_{13}$ plus $A_3$; $T_2+A_1$; $I+T_2$; or $O_7+A_1$&
 Respectively: $53/70+1/4$; $53/70+1/7+1/10=1$;
 X7; or $233/300+9/100+2/15=1$\\
10&$A_{14}$ plus $T_2$; $2A_1$; $I+A_1$; $N_7$; or $O_7$&
 Respectively: $16/19+3/19=1$; $4/5+.2=1$;
 $4/5+.20204$; or $.91286445198+.09$ for either typed branch\\
11&$A_{15}$ plus ordinary $A_1$ or actual $I$&X8 in its
 complete isolated/nonisolated rank-one domains\\
12&Entire $A_{16}$&X9\\
13&$B_{14}+N_7$ or $B_{14}+O_7$&
 $.91859+.09=1.00859$\\
14&$B_{15}$ plus ordinary $A_1$ or actual $I$&
 X10 for ordinary $A_1$; X11 for actual $I$\\
15&Entire $B_{16}$&$e_{B_{16}}\ge1.00640742675$, after
 its own domain inclusion and edge estimate for the unique survivor\\
16&$D_{13}$ plus $I+T_2$ or $O_7+A_1$&
 $.82436$ plus the remainder floor $.21$; joint on $I+T_2$,
 and coarsened from $.09+2/15$ on $O_7+A_1$\\
17&$D_{14}$ plus ordinary $I_2(m)$, $m\ge4$; $N_7$; or $O_7$&
 Ordinary: $.84816+3/19$. Either typed branch:
 $.91859+.09$\\
18&$D_{15}$ plus ordinary $A_1$ or actual $I$&
 X10 for ordinary $A_1$; X12 for actual $I$\\
19&Entire $D_{16}$&X13, the actual quadratic reduction and
 complete dyadic argument\\\bottomrule
\end{longtable}

Each sum is at least one, which suffices against the strict budget
\eqref{c:budget}; no uniform positive excess is needed. For example,
group 5 uses three disjoint rows summing to exactly one. Group 4
instead uses the joint support $A_{11}+I$:
\[
 e_{A_{11}}+e_I\ge18/25,\qquad e_{T_4}\ge2/7.
\]
These rows have disjoint supports. The floor $18/25$ must not be
assigned to $A_{11}$ alone or combined with a second full use of $I$.

The nineteen groups give $49$ branches. Together the two covers have
$88$: $72$ energy branches and $16$ local contradiction branches,
with unbounded dihedral parameters quantified within their full domains.

\subsection{The arithmetic local contradiction statements}
\label{a:local-exit-section}

For completeness, Table~\ref{a:local-exits} specifies every local
contradiction used by Tables~\ref{a:rank17} and~\ref{a:rank16}. These
are imported local theorems with their full arithmetic and geometric
hypotheses. They are not extra scalar energy claims. Their LP row is
zero mass on the exactly matched domain.

\begin{longtable}{>{\raggedright\arraybackslash}p{.07\linewidth}
 >{\raggedright\arraybackslash}p{.31\linewidth}
 >{\raggedright\arraybackslash}p{.52\linewidth}}
\caption{Arithmetic local contradiction inputs. ``Entire'' refers to
the complete original vertex forest.}\label{a:local-exits}\\
\toprule
Code&Matched entire forest&Complete local argument retained\\\midrule
\endfirsthead
\toprule Code&Matched entire forest&Local argument\\\midrule\endhead
X1&$A_{15}+T_2$, rank $17$&The ordinary-dihedral theorem,
 including every silver, golden and low-energy label domain.\\
X2&$A_{16}+A_1$ or $A_{16}+I$, rank $17$&The complete
 actual rank-one-opposite theorem, isolated and nonisolated.\\
X3&Entire $A_{17}$, rank $17$&The endpoint and field reductions,
 with proved finite-anchor inclusions in outputs (i) and (ii).\\
X4&$B_{16}+I$, rank $17$&The complete six-role short-cap theorem,
 with its quadratic and cubic field alternatives.\\
X5&$D_{16}+A_1$, rank $17$, with ordinary $A_1$&The ordinary
 rank-one theorem, retaining both fork alternatives.\\
X6&$D_{16}+I$, rank $17$&The actual root-line theorem in degrees
 two, three and four, including both signed forks and all dyadic cases.\\
X7&$A_{13}+I+T_2$, rank $16$&The joint endpoint-and-remainder
 theorem with the complete ordinary dihedral domain.\\
X8&$A_{15}+A_1$ or $A_{15}+I$, rank $16$&The complete
 actual rank-one-opposite theorem, with all height boundaries.\\
X9&Entire $A_{16}$, rank $16$&The endpoint, field, finite-anchor
 and time-ratio theorem for this whole component.\\
X10&$B_{15}+A_1$ or $D_{15}+A_1$, rank $16$, with ordinary $A_1$&
 Their respective nonisolated reductions and ordinary isolated theorem,
 including $h=1/2$.\\
X11&$B_{15}+I$, rank $16$&The complete isolated-short theorem
 for $0<h<1/2$.\\
X12&$D_{15}+I$, rank $16$&The complete isolated-short theorem
 for $0<h<1/2$, including both signed forks.\\
X13&Entire $D_{16}$, rank $16$&The quadratic reduction and manual
 dyadic theorem: odd, inert, ramified and split cases.\\\bottomrule
\end{longtable}

The rank-one theorems follow the same actual opposite wall. Nonzero
contact with the large component gives the nonisolated alternative;
otherwise its actual isolated height lies in $0<h<1/2$ or $h\ge1/2$.
The equality case belongs to the ordinary theorem. These alternatives
cover an existing wall without replacing it by a formal extension.

The whole-$B_{17}$ and whole-$D_{17}$ energy rows use separate proved
inclusions in output (iii) and separate energy estimates for its two
survivors. Output (iv) likewise retains its unique pair until the actual
$B_{16}$ edge estimate is applied. These outputs are not emptiness
claims; whole $D_{16}$ instead uses the manual theorem X13.

\subsection{One joined arithmetic certificate}
\label{a:certificate}

The arithmetic tables give a hand certificate. On each energy branch in
Tables~\ref{a:rank17} and~\ref{a:rank16}, give the displayed disjoint
support rows weight one and all additional rows weight zero. Each
original energy has total weight at most one, and the weighted lower
bounds sum to at least one. A branch X1--X13 has zero mass and hence
zero normalized deficit by its retained local theorem.

The actual computation uses the same assembler for all $88$ arithmetic
branches and for the $829$ compact branches at dimensions $18$--$29$.
The joined catalogue has $917$ branches, including $167$ local
contradiction domains; its checked capacity dual has $3324$ variables
and $2488$ rows. The reproduction section records these global LP
checks separately from the retained local proofs.

\begin{theorem}[Compact arithmetic recovery]\label{a:recovery}
Retain Input~\ref{a:input-package}, the compact local inputs of
Section~\ref{c:section}, and the published compact ceiling $29$.
The joined energy LP is infeasible for every actual dimension at least
$16$. It recovers the existing compact arithmetic bound $n\le15$.
\end{theorem}
\begin{proof}
Every actual arithmetic configuration at ranks $16$ and $17$ belongs
to the complete corresponding cover. A branch with a retained local
contradiction is impossible. On every other branch the exact capacity
certificate gives total original energy at least one, contrary to
\eqref{c:budget}. The compact domains at ranks $18$--$29$ apply unchanged,
since an arithmetic compact polytope is compact; their certificate was
proved independently of arithmetic in Theorem~\ref{c:recovery}.
The external ceiling excludes dimensions $30$ and above.
\end{proof}



\section{Coverage of the dimension thresholds}
\label{sec:thresholds}

Proposition~\ref{prop:threshold} is the monotonicity argument for all
three cases. What differs is the proof that the fixed input family
covers every actual dimension in the excluded range:

\begin{center}\small
\begin{tabular}{>{\raggedright\arraybackslash}p{.17\linewidth}
 >{\raggedright\arraybackslash}p{.35\linewidth}
 >{\raggedright\arraybackslash}p{.38\linewidth}}
\toprule
Case&Certificate domains&Coverage beyond the finite endpoint\\\midrule
Finite volume&Finite and ideal endpoint at $n=58$&Analytic face and
 vertex estimates, with the stated ideal descent, cover every $n\ge59$\\
Compact&Complete catalogue at $18\le n\le29$&Published compact
 ceiling $n\le29$\\
Compact arithmetic&Arithmetic covers at $n=16,17$, joined to the
 same compact catalogue at $18\le n\le29$&The same published compact
 ceiling\\\bottomrule
\end{tabular}
\end{center}

Vinberg's ceiling excludes compact polytopes in dimensions
$n\ge30$~\cite{Vinberg1984,Vinberg1985}. It is independent input.
The arithmetic formulation copies the compact domains unchanged;
it does not insert the conclusion $n\le17$ as a row. The
finite-volume tail is part of its geometric input and certificate,
not an assertion that a new finite matrix was exported at every
higher dimension.

Thus the exclusions concern every dimension above the respective
threshold, rather than three isolated exact dimensions. Passing to
an arbitrary face or deleting diagram nodes is never needed.
To lower a threshold one would have to supply complete valid data in
the newly admitted dimensions; changing $N$ alone supplies none.

\section{Exact certificates and the interpretation of the computations}
\label{sec:certificates}

The numerical part of the two compact recoveries is small after the
local theorems have been stated. Its purpose is to find the global
coefficients. It does not reconstruct the geometric proof of each
local estimate. The matrix stores the support of every estimate,
its rational lower bound, and the domain on which it is used.

The capacity constraints are assembled for all nonexcluded domains
at once. A numerical solver proposes their multipliers. Those
multipliers are converted to rational numbers and every matrix row
and variable bound is checked exactly. A second run regenerates the
entire configuration catalogue and both the normalized primal and
capacity matrices, compares their complete bytes, and rechecks the
rational point. Thus the final implication depends on a rational
inequality check, not on floating-point tolerances or a solver's
status message.

\begin{center}\small
\begin{tabular}{lrrrr}
\toprule
Problem&Domains&\shortstack{Local\\exclusions}&
\shortstack{Primal\\variables/rows}&\shortstack{Dual\\variables/rows}\\
\midrule
Compact&829&151&$3471/4876$&$3157/2224$\\
\shortstack[l]{Arithmetic,\\with compact tail}&917&167&$3869/5147$&$3324/2488$\\
\bottomrule
\end{tabular}
\end{center}
The finite-volume computation has a different representation. Its
local matrix has $5\,810$ variables and $103\,606$ rows. The complete
finite-$58$ network certificate is assembled from verified blocks
sharing one coefficient assignment. The exact feasible certificate
and the eight-row minimax lower certificate match at $377$. No
numerical optimizer was run for that result, and no claim is made
that its entire coupled matrix was flattened into one file. A
subsequent check repeated the specified local, coverage-assembly
and lower-bound tasks while reusing the original network executions.

Some local inequalities are already dominated by others on the same
domain and support. For example, the prefix estimate with the extra
initial gain is stronger than the corresponding prefix estimate
alone. Keeping both is harmless and makes their sources visible.
More generally, a zero multiplier in one certificate only says that
that certificate did not use the row. It does not prove that the row
is logically redundant. We do not claim a minimal list of geometric
inputs.

\section{Scope for further constraints}
\label{sec:strengthening}

A new valid condition may shrink an admitted domain or add a row to
its cone. A larger observable or multiplier family may also enlarge
the available certificates. In each case the actual-object lift and
the fixed dimension labels must be preserved. These are the
requirements for reusing the same linear deduction and monotonicity.

The positive-definiteness, Schur-complement and arithmetic arguments
currently enter through local theorems. Replacing some by explicit
matrix inequalities could give an SDP. Real Gram positivity alone
would not replace integral Cartan cycles, nonphysical embeddings or
the stated lattice-group inclusion. Such a replacement must say which
local facts it proves and which remain inputs. The present results
concern the exact LP recovery of~\eqref{eq:three-bounds}.

\appendix
\section{Source statements, data and reproduction}
\label{app:reproduction}

The mathematical sources are the dimension-bound article and its
complete supplement~\cite{BoundsArticle,BoundsSupplement}. The
finite-volume LP also uses the common-law transfers and endpoint/tail
composition described in the research notes. The associated repository
\cite{LPRepository} contains the reading guide
\path{notes/lp/UNIFORM_LP_PARADIGM.md} and the detailed result note
\path{notes/lp/existing_proof_lp_recovery_results_20261118.md}.
These distinguish the mathematical premises, source reviews,
actual computations and exact certificate checks.

The following paths identify the precise records in the associated
research checkout. They are not additional mathematical hypotheses.
The repository's reproduction guide distinguishes the files supplied
with the public checkout from the separate bulk research inputs.

\begingroup\small\raggedright
\begin{longtable}{>{\raggedright\arraybackslash}p{.23\linewidth}
 >{\raggedright\arraybackslash}p{.70\linewidth}}
\toprule
Record&Path\\\midrule\endfirsthead
\toprule Record&Path\\\midrule\endhead
Finite-volume completion&\path{audit/joint58_complete_lp_execution_20261110.json}\\
Later declared rerun&\path{audit/lp_rerun_report_execution_20261115.json}\\
Compact local-data rules&\path{scripts/compact17_proof_recovery_energy_catalogue_20261118.py}\\
Arithmetic local-data rules&\path{scripts/existing_proof_arithmetic_energy_domains_20261118.py}\\
Common energy assembler&\path{scripts/run_existing_proof_energy_lp_20261118.py}\\
Exact matrix transport/check&\path{scripts/finite_volume_lp_matrix_20261106.py}\\
Compact complete models&\path{results/existing_proof_compact17_energy_master01_20261118/}\\
Arithmetic complete models&\path{results/existing_proof_arithmetic15_energy_master01_20261118/}\\
Independent execution review&\path{audit/existing_proof_energy_lp_execution_independent_review_20261118.json}\\
Final recovery acceptance&\path{audit/existing_proof_lp_recovery_completion_root_acceptance_20261118.json}\\
\bottomrule
\end{longtable}
\endgroup

Each compact model directory contains the complete configuration
catalogue, normalized primal, capacity dual, rational candidate,
construction report and fresh verification report. The monitor
records give the exact commands, source hashes, resource samples,
elapsed times and memory measurements. Source reviews check the
configuration-to-theorem applicability; the exact checker verifies
the final rational linear algebra. A checksum establishes byte
identity, and does not replace either of those arguments.

For both compact cases, reproduction uses the existing bounded runner
with its shared exclusive lock. The construction stage is
\texttt{export-solve}; the separate stage \texttt{verify} regenerates
both matrices and rechecks the candidate. Arithmetic supplies both
catalogue modules. Output directories and monitor ledgers must be
fresh. The recorded runs use a $59$-second, $1$-GiB, one-thread outer
watchdog and retain the original percentage resource policy.

The computational statement is limited to the specified models.
The local proofs, complete source domains, original field reductions,
four arithmetic computed-output premises and applicable lattice
theorems retain their stated roles. The successful local scientific
checks and formal modules were not reexecuted merely to construct
these global LPs. Typesetting this manuscript performs no
mathematical verification.

\paragraph{Acknowledgments.}
The author used OpenAI's Codex as a research and coding assistant to
investigate proof ideas, develop the linear programming formulations,
improve the exposition, and generate and audit the accompanying Python
code. The author reviewed the resulting mathematics and code and takes
responsibility for the final content.

\begin{thebibliography}{99}
\bibitem{GuoZheng2026}
L.~Guo and F.~Zheng, \emph{Dimension bounds for finite-volume hyperbolic
Coxeter polytopes}, preprint, arXiv:2611.01795v1, 2026.
\url{https://arxiv.org/abs/2611.01795}.

\bibitem{BoundsArticle}
A.~Kolpakov, \emph{Dimension bounds for hyperbolic Coxeter polytopes},
working article, 2026.
\url{https://github.com/sashakolpakov/CoxeterPolytopes/tree/main/arxiv/hyperbolic_coxeter_dimension_bounds}.

\bibitem{BoundsSupplement}
A.~Kolpakov, \emph{Dimension bounds for hyperbolic Coxeter polytopes:
supplementary proofs}, 2026.
\url{https://github.com/sashakolpakov/CoxeterPolytopes/tree/main/manuscripts/addenda/dimension_bounds_20261030}.

\bibitem{LPRepository}
A.~Kolpakov, \emph{CoxeterPolytopes: proof sources and LP research
records}, research repository, 2026.
\url{https://github.com/sashakolpakov/CoxeterPolytopes}.

\bibitem{Vinberg1984}
E.~B.~Vinberg, The absence of crystallographic groups of reflections
in Lobachevsky spaces of large dimension,
\emph{Trudy Moskov. Mat. Obshch.} \textbf{47} (1984), 68--102;
English translation, \emph{Trans. Moscow Math. Soc.} (1985), 75--112.
\url{https://www.mathnet.ru/eng/mmo444}.

\bibitem{Vinberg1985}
E.~B.~Vinberg, Hyperbolic reflection groups,
\emph{Russian Math. Surveys} \textbf{40}:1 (1985), 31--75.
\url{https://doi.org/10.1070/RM1985v040n01ABEH003527}.

\bibitem{Bugaenko1984}
V.~O.~Bugaenko, Groups of automorphisms of unimodular hyperbolic
quadratic forms over $\Z[(\sqrt5+1)/2]$,
\emph{Vestnik Moskov. Univ. Ser.~1 Mat. Mekh.} (1984), no.~5, 6--12.
\url{https://www.mathnet.ru/eng/vmumm3408}.
\end{thebibliography}
\end{document}
